\documentclass[11pt,letterpaper]{article}
\usepackage[T1]{fontenc}
\usepackage[utf8]{inputenc}
\usepackage{lmodern,microtype}
\usepackage[margin=0.92in]{geometry}
\usepackage{amsmath,amssymb,amsthm,mathtools}
\usepackage{booktabs,longtable,array,enumitem}
\usepackage[dvipsnames]{xcolor}
\usepackage{listings}
\usepackage{xurl}
\usepackage{hyperref}
\usepackage{fancyhdr}
\hypersetup{colorlinks=true,linkcolor=MidnightBlue,citecolor=MidnightBlue,
 urlcolor=MidnightBlue,pdftitle={Finite observations and uniform constructions in iPRA},
 pdfauthor={Prepared for Vladimir Reshetnikov},bookmarksnumbered=true}
\urlstyle{same}
\pagestyle{fancy}\fancyhf{}
\fancyhead[L]{\small Finite observations in iPRA}
\fancyhead[R]{\small Research framework and applications}
\fancyfoot[C]{\thepage}
\setlength{\headheight}{14pt}
\setlength{\emergencystretch}{3em}
\setlist{topsep=5pt,itemsep=3pt,parsep=2pt}
\lstset{basicstyle=\small\ttfamily,columns=fullflexible,breaklines=true,
 frame=single,rulecolor=\color{black!20},showstringspaces=false}
\newcommand{\iPRA}{\ensuremath{\mathsf{iPRA}}}
\newcommand{\N}{\mathbb N}
\newcommand{\Q}{\mathbb Q}
\newcommand{\PR}{\mathrm{PR}}
\newtheorem{theorem}{Theorem}[section]
\newtheorem{lemma}[theorem]{Lemma}
\newtheorem{proposition}[theorem]{Proposition}
\newtheorem{corollary}[theorem]{Corollary}
\theoremstyle{definition}\newtheorem{definition}[theorem]{Definition}
\theoremstyle{remark}\newtheorem{remark}[theorem]{Remark}
\title{\textbf{Finite observations and uniform\\constructions in iPRA}\\[7pt]
\large A compositional framework, degree-uniform factorization,\\
convexity moduli, and finite algebraic sections}
\author{Prepared for Vladimir Reshetnikov}
\date{November 15, 2026}
\begin{document}
\maketitle
\begin{abstract}
We develop a common method for translating presented mathematical objects
and quantitative proofs into intuitionistic primitive recursive arithmetic.
The central construction replaces a variable-depth chain of operations on
infinite names by backward numerical demand bounds and forward finite
tables. Evaluation demands and verification demands are kept distinct;
finite Boolean invariants justify the permitted induction.

The framework yields three concrete advances relative to the inspected
ProveIt reports. First, using the existing finite rational descent lemma,
we obtain one full-factorization program uniform in the numerical degree
for each fixed primitive recursive coefficient evaluator. Second,
positive rational multiplicative convexity supplies the candidate's local
continuity modulus in the Bohr--Mollerup argument, from two coarse
endpoint evaluations. Third, for $y^r=f(x)$ with $m=\deg f>rs$, every
proper linear subspace contains feature images of at most $m$ distinct
rational points under $(1,x,\ldots,x^s,y)$; this supplies an explicit
zero-count input to a finite-cover route toward Diophantine finiteness.

Additional theorems construct infimum names from finite nets and coherent
minimizers from quantitative uniqueness, including a complete strong
convexity instance. Exact-input, apartness, uniformity, and ineffective-bound
obstructions are retained explicitly. The archive includes a checked Lean
finite-table kernel, type-checked mathematical target specifications,
exact arithmetic examples, and a detailed path to deeply embedded
\texttt{PRAAx} derivations. It does not claim that those full object-theory
derivations have already been completed.
\end{abstract}
\newpage
\tableofcontents
\newpage
\section{Results, scope, and the repository baseline}
\label{sec:overview}

The useful common framework is not a blanket principle saying that
computable mathematics belongs to \iPRA. It is a disciplined interface
between finite observations, numerical bounds, and proofs. Many of the
repository's reports already use parts of this interface. We consolidate
those parts into compositional theorems, strengthen the interface where
nested constructions need it, and use it to prove specific new extensions.

The inspected repository revision is
\begin{center}
\texttt{94000996e69ab0c527393a1d02a4a037a4eb0493}.
\end{center}
The directory inventory contains 48 delivered proof packets. Its own
README explicitly says that placement did not rerun or validate all the
packets, and that many Lean files are outlines \cite{repo-index}. Our
audit likewise does not certify every theorem in that collection. The
provenance ledger records the particular files used, their repository
blob hashes, and the source locations of the obstacles addressed below.

\subsection{What is established here}

\begin{longtable}{@{}>{\raggedright\arraybackslash}p{0.24\linewidth}>{\raggedright\arraybackslash}p{0.46\linewidth}>{\raggedright\arraybackslash}p{0.24\linewidth}@{}}
\toprule
Result & Mathematical content & Dependence and status\\
\midrule
\endhead
Finite observation compilation & Variable-depth operations are evaluated
by a bounded input table, a backward demand schedule, and a forward
numerical recursion. Finite contracts propagate by quantifier-free
induction. & General human proof; selected locality and compilation
lemmas checked in Lean.\\[5pt]
Uniform polynomial factorization & For each fixed PR coefficient evaluator,
one PR program factors monic polynomials uniformly in numerical degree,
returning coherent complex names and explicit coefficient-error bounds.
& Uses the existing finite rational descent lemma; all new normalization,
deflation, and iteration estimates proved here.\\[5pt]
Derived Bohr--Mollerup modulus & Positive rational multiplicative convexity
gives a local Lipschitz bound from two endpoint samples, so the candidate
modulus is no longer independent input. & Direct elementary proof;
canonical extension of the rational presentation.\\[5pt]
Finite-net extrema and coherent minimizers & A net and a supplied modulus
give an infimum name; quantitative uniqueness gives a coherent optimizer.
Strong convexity provides an explicit uniqueness modulus. & Direct
finite constructions and complete quantitative proofs.\\[5pt]
Superelliptic section bound & For $y^r=f(x)$, $\deg f=m>rs$, a proper
subspace meets $(1,x,\ldots,x^s,y)$ in at most $m$ distinct point images.
& Unconditional finite algebra. A tuple-local $M$-subspace cover then
gives the conditional cardinality bound $mM$.\\
\bottomrule
\end{longtable}

The new factorization theorem closes the fixed-degree restriction in the
existing FTA discussion, with its coefficient evaluator still fixed in the
correct first-order sense. The new convexity argument answers an explicit
further-research question in the Bohr--Mollerup report. The Diophantine
result discharges one geometric input for a defined family; it does not
supply the quantitative approximation theorem that must produce the
covers.

\subsection{The central organization}

The following separation makes mathematical reuse safe.
\begin{enumerate}
\item \textbf{Representation.} State exactly what finite observations
      are available. Exact algebraic data, Cauchy names, formal
      coefficients, and names with apartness certificates are distinct
      interfaces.
\item \textbf{Construction.} Give primitive recursive numerical
      procedures with an explicit demand on their input observations.
\item \textbf{Verification.} State which finitely many input contracts
      establish the requested output facts, and bound their demand
      separately from the construction's own queries.
\item \textbf{Arithmetic proof.} Fix the numerical output challenge,
      compile all its verification demands, and induct on a Boolean
      property of the resulting finite data.
\item \textbf{Interpretation.} Only after this finite argument quantify
      over the arbitrary output challenge and read the result as an
      identity, a continuous object, a root, or a finiteness assertion.
\end{enumerate}

This organization is useful even when it does not remove a hypothesis.
For example, it pinpoints the analytic estimate still needed for a
stronger prime-distribution theorem, instead of hiding that estimate
inside an interface called ``effective asymptotics.'' It also explains
why a finite cardinality theorem does not automatically enumerate all
solutions.

\subsection{How to read the report}

Sections~\ref{sec:foundations} and~\ref{sec:compiler} give the shared
arithmetic language and the compilation theorem.
Section~\ref{sec:fta-extension} is the most detailed variable-depth
application. Section~\ref{sec:quantitative-analysis} treats precision,
extrema, uniqueness, and convexity. Section~\ref{sec:finite-sections}
provides the elementary algebraic zero count and its conditional
Diophantine interface. The final sections explain exactly which reported
obstacles remain, how to complete the Lean development, and what was
actually checked in this archive.

\section{The arithmetic boundary and the representation discipline}
\label{sec:foundations}

\subsection{The precise base theory}

Throughout, \iPRA{} means intuitionistic first-order arithmetic in the
language containing a symbol for every primitive recursive numerical
function, with their defining equations, successor injectivity and
separation, and induction for quantifier-free formulas. All variables in
an internal theorem range over natural numbers. The repository implements
exactly this theory as \texttt{PRAAx}; its primitive recursive symbols are
\texttt{PR k}, its arithmetic terms are \texttt{PTm n}, and its derivations
are \texttt{Prf PRAAx}. The native scalar presentation is connected to it by
\texttt{pra\_iff\_native} \cite{repo-pr,repo-native}.

An assertion about a primitive recursive presentation $A$ is a
\emph{schema}: fix its finite primitive recursive definition externally,
then substitute that definition into a first-order formula. Numerical
parameters of $A$, including degree, dimension, precision, and a finite
input polynomial, remain ordinary quantified variables. We will never
infer that a total computable function is primitive recursive, nor that a
primitive recursive predicate has its universal truth provable in
\iPRA. The latter would confuse representation with proof.

We use natural codes for signed integers, rational numbers with positive
denominator, finite tuples, and finite lists. One explicit list coding
suffices: let
\[
 \langle a,b\rangle=(a+b)^2+a+1,
 \qquad []=0,\qquad [a;L]=\langle a,L\rangle.
\]
Both components of a nonzero pair are less than its code. Recover them
by a search bounded by the code, with a default value on malformed
inputs. Injectivity follows by separating successive squares and then
subtracting. A traversal can be bounded by the initial code; an explicitly
stored length permits a sharper bound. All accesses, finite maxima,
appends, rational cross-multiplication tests, and their totalized default
branches are primitive recursive. These facts concern the numerical
algorithms; their algebraic laws follow by the corresponding bounded
inductions. Nothing depends on efficient binary encodings.

\begin{lemma}[Boolean reflection and bounded induction]
\label{lem:reflection}
For each bounded first-order formula $B(\mathbf x)$ built from primitive
recursive predicates there is a primitive recursive Boolean
$\chi_B(\mathbf x)\in\{0,1\}$ such that
\[
 \iPRA\vdash B(\mathbf x)\leftrightarrow \chi_B(\mathbf x)=1.
\]
Consequently induction on $B$ is derivable using quantifier-free induction.
The formula $B$ is fixed externally; this is not a uniform truth predicate
for all bounded formulas with arbitrary primitive recursive symbols.
\end{lemma}
\begin{proof}
Equality and numerical order have primitive recursive characteristic
functions. For example compare two numbers by simultaneous truncated
subtraction. Their defining recursions prove decidability by
quantifier-free induction. Use Boolean complement, conjunction, and
disjunction for the corresponding connectives. Since the already
reflected subformulas are decidable, implication is represented by
$\max(1-\chi_C,\chi_D)$; this use of Boolean implication makes no assertion
of excluded middle for arbitrary formulas.

A bounded universal quantifier is represented by the fold
\[
 F(0)=1,\qquad F(j+1)=F(j)\chi_C(j,\mathbf x).
\]
Here is a noncircular proof of its reflection law. With $i$ free,
quantifier-free induction on $j$ proves
$i<j\wedge F(j)=1\longrightarrow\chi_C(i,\mathbf x)=1$; the successor
step uses the cases $i<j$ and $i=j$. Conversely fix a final $n$ and
keep $\forall i<n\ \chi_C(i,\mathbf x)=1$ as a context assumption.
Induction on the quantifier-free formula $j\le n\longrightarrow F(j)=1$
proves $F(n)=1$. Thus bounded induction is not presupposed in this
proof of bounded reflection. The existential fold uses
$1-\prod(1-\chi_C)$. Its witness direction can also use the bounded
first-success scan described below. Formula induction here is an external
induction on a fixed syntax tree; only the folds' numerical inductions
occur inside \iPRA. Finally replace $B$ in its base and step by the
provably equivalent equation $\chi_B=1$, apply the permitted induction,
and replace back.
\end{proof}

\begin{lemma}[Bounded witness extraction]
\label{lem:bounded-choice}
Let $R(\mathbf x,i,y)$ be a fixed primitive recursive relation and
$b(\mathbf x,i)$ a fixed primitive recursive bound. There is a primitive
recursive $w(\mathbf x,i)$ such that \iPRA{} proves
\[
 (\exists y\le b(\mathbf x,i)\ R(\mathbf x,i,y))
 \ \longrightarrow\ R(\mathbf x,i,w(\mathbf x,i))
       \ \wedge\ w(\mathbf x,i)\le b(\mathbf x,i).
\]
For any numerical $m$, the list of these witnesses for $i<m$ is primitive
recursive and satisfies every required relation whenever the $m$ bounded
existences hold.
\end{lemma}
\begin{proof}
Scan $0,\ldots,b$, keeping a flag indicating whether a successful test
has occurred and, if so, its first index. The invariant says that either
no scanned test succeeds, or the saved index is in the scanned range,
satisfies $R$, and precedes every other success. All its quantifiers have
bounds supplied by the scan counter. Lemma~\ref{lem:reflection} permits
induction on this invariant. A supplied success forces the second case
at the end. In the no-success case return zero; this makes the operation
total. A second bounded fold appends the $m$ outputs. This is a construction
of a list, not an invocation of a choice axiom.
\end{proof}

\subsection{What a presented infinite object is}

A \emph{presentation} has three parts:
\begin{enumerate}
\item a fixed primitive recursive observation function
      $A(\mathbf a,j)$ with numerical parameters $\mathbf a$;
\item finite metadata, such as an amplitude bound, a dimension, a
      separation witness, or a primitive recursively supplied modulus;
\item arithmetic contracts on its observations. Frequently these have the
      form $\forall\mathbf u\ R(\mathbf a,\mathbf u)$ for a primitive
      recursive relation $R$.
\end{enumerate}
Examples include a finite matrix, an exact algebraic presentation, a
rational approximation to a complex number, the coefficients of a formal
power series, a finite probability law, and a rational sampling routine
for a function. A fixed function can have several parameters; that does
not turn it into a universal evaluator of arbitrary programs.

Different contracts express different mathematical content. A formal
power series needs no convergence modulus to state identities of finitely
many coefficients. Evaluation at a point additionally needs a quantitative
tail estimate. A Cauchy name permits approximate equality tests, but an
exact algebraic presentation can provide a decidable zero test. The
framework keeps these interfaces distinct because algorithms really do
use different information.

\begin{definition}[Regular rational names]
A fixed rational observation function $a(k)$ is regular when
\begin{equation}
 |a(k)-a(\ell)|\le 2^{-k}+2^{-\ell}
 \qquad(k,\ell\in\N).                         \tag{RN}\label{eq:RN}
\end{equation}
Complex names use pairs of rational observations and the rational
squared-norm version of the same estimate. We write $a\sim b$ when
\[
 \forall k\quad |a(k)-b(k)|\le 2^{1-k}.
\]
This is an arithmetic relation on two fixed presentations, not a
quotient type assumed as an object of first-order arithmetic. Its
usual real-number reading is explanatory notation.
\end{definition}

\begin{lemma}[Eliminating an arbitrary precision error]
\label{lem:vanishing-error}
For rational $q,c$ and nonnegative rational $C$, if
$q\le c+C2^{-r}$ for every $r$, then $q\le c$.
This implication is provable in \iPRA{} as a schema and requires neither
Markov's principle nor an unbounded minimization.
\end{lemma}
\begin{proof}
Compare the rational numbers $q,c$. If $q\le c$, finish. Otherwise write
$q-c=v/w$ with positive integers $v,w$, and write $C=a/b$ with $b\ge1$.
Choose, for example, $r=aw+2$. Then $2^r>aw$ (including $a=0$), so
$C2^{-r}<1/w\le v/w$. Substitution of this particular $r$ contradicts
the assumed inequality. The exponential inequality follows by ordinary
quantifier-free induction. The case distinction is rational and decidable.
\end{proof}

This lemma supplies elementary passages to exact name equality. For
example, transitivity of $\sim$ follows by comparing $a(k)$ and $c(k)$
through $a(r),b(r),c(r)$: regularity and the two equivalences give
\[
 |a(k)-c(k)|\le 2^{1-k}+6\,2^{-r}.
\]
Eliminate the last term by Lemma~\ref{lem:vanishing-error}. Thus no real
completion theorem is required for this equivalence law.

\begin{proposition}[Arithmetic on names with budgets]
\label{prop:name-arithmetic}
Regular names admit primitive recursive addition and multiplication with
explicit precision shifts. Inversion admits such a construction when a
positive apartness certificate is supplied.
\end{proposition}
\begin{proof}
For addition use $c(k)=a(k+2)+b(k+2)$ and add the two regularity bounds.
For multiplication suppose every sample of each input has absolute value
at most the supplied integer $M\ge1$. Choose $s$ with $2^s\ge2M$ and set
$c(k)=a(k+s)b(k+s)$. The identity
\[
 ab-a'b'=a(b-b')+b'(a-a')
\]
bounds the difference by $2M(2^{-k-s}+2^{-\ell-s})$, as required.
A cap can itself be derived from a zero-th sample and regularity:
$|a(k)|\le|a(0)|+2$.

For inversion suppose $|a(h+2)|\ge2^{-h}$. Every sample of index
$j\ge h+2$ then has $|a(j)|\ge2^{-h-1}$. Put $s=2h+3$ and
$c(k)=1/a(k+s)$, with a default zero value if a denominator test fails.
On valid inputs it does not fail, and
\[
 |c(k)-c(\ell)|\le
   2^{2h+2}(2^{-k-s}+2^{-\ell-s})
   \le2^{-k}+2^{-\ell}.
\]
The same finite difference identities show extensionality under $\sim$;
any remaining error is removed by Lemma~\ref{lem:vanishing-error}.
All regularity proofs fix the two requested precisions before any
induction on a finite sum or product.
\end{proof}

The apartness witness is substantive. The premise $\neg(a\sim0)$ does
not itself specify one. Likewise, a root-producing procedure may depend
on the presentation and choose different valid roots from equivalent
coefficient names. We demand extensionality of the asserted mathematical
relation, not an impossible global continuous choice of a preferred root.

\subsection{Three claims that must not be conflated}

We use the following status vocabulary throughout the report.
\begin{description}[style=nextline]
\item[Mathematical derivation in the stated arithmetic.]
A human-readable argument with primitive recursive constructions and
identified quantifier-free induction invariants. This is the main level
of the present contribution.
\item[Ambient Lean verification.]
A theorem checked by Lean about finite tables, streams, or arithmetic.
This checks that theorem in Lean's foundations. It does not by itself
produce a proof in the much weaker object theory.
\item[Deep iPRA verification.]
An actual term of type $\Gamma\mathbin{\vdash}[\mathtt{PRAAx}]\varphi$,
with the primitive recursive definitions and the formula $\varphi$
explicitly encoded. The repository already provides the logical backend;
the new mathematical compilation layer is not yet fully formalized.
\end{description}

The distinction is constructive, not merely terminological: a Lean proof
by induction on an arbitrary proposition does not authorize that induction
inside \iPRA. The next section gives the finite invariant that makes
specific uses legitimate.

\section{A compositional finite-observation framework}
\label{sec:compiler}

\subsection{Finite safety under infinite hypotheses}

A frequent mistake is to justify an induction by saying that every
individual computation is finite. The induction formula itself must be
permitted. The following lemma gives the needed shape.

\begin{theorem}[Finite safety compilation]
\label{thm:safety}
Fix primitive recursive functions $\operatorname{step}(\mathbf x,j,s)$,
$s_0(\mathbf x)$, a horizon $H(\mathbf x)$, and a Boolean invariant
$I(\mathbf x,j,s)$. Let $\Gamma(\mathbf x)$ be any fixed collection of
first-order assumptions. Suppose there are \iPRA{} derivations, under
$\Gamma$, of
\[
 I(\mathbf x,0,s_0)=1,
 \qquad
 j<H\ \wedge\ I(\mathbf x,j,s)=1
 \ \longrightarrow\ I(\mathbf x,j+1,\operatorname{step}(\mathbf x,j,s))=1.
\]
Define $s_{j+1}=\operatorname{step}(\mathbf x,j,s_j)$ by primitive
recursion. Then there is an \iPRA{} derivation under $\Gamma$ of
\[
 \forall j\le H\quad I(\mathbf x,j,s_j)=1.
\]
This is a proof-transformation theorem: the base and step derivations
are required, not inferred merely from their truth in $\N$.
\end{theorem}
\begin{proof}
Keep $\mathbf x,H$ fixed. Induct on the quantifier-free formula
\[
 j\le H\ \longrightarrow\ I(\mathbf x,j,s_j)=1.
\]
At a successor, a rational/numerical case split first checks $j+1\le H$.
If it fails the implication follows. If it holds, $j<H$ and $j\le H$;
the induction hypothesis and the supplied step derivation give the
conclusion. The hypotheses $\Gamma$ remain in the context. Quantifier-free
induction does not require the other context formulas to be quantifier-free.
Ordinary universal introduction subsequently closes the numerical
parameters. Syntactically, encode this guarded formula as a
\texttt{Fm PTm}; its \texttt{QF} proof permits the single application
of \texttt{PRA.qfInd}.
\end{proof}

A finite contract can accumulate many facts: bounds on a state,
well-formedness of a list, every earlier local transition check,
a polynomial coefficient identity, and every relevant approximation
inequality. A bounded conjunction of those facts becomes one Boolean
invariant by Lemma~\ref{lem:reflection}. A universal name hypothesis is
used only to verify the finitely many input observations required by that
particular run.

\subsection{Two different budgets}

An observation procedure has two potentially different resource bounds:
\begin{enumerate}
\item its \emph{evaluation demand} is the largest input observation needed
      to compute an output observation;
\item its \emph{verification demand} is the largest input observation and
      contract instance needed to prove the desired finite output facts.
\end{enumerate}
The second can be larger. A root approximation at precision $k$ might
query one finite prefix, whereas comparing it to another approximation or
proving a residual inequality asks for additional rows. An implementation
trace alone does not bound its mathematical proof obligations. We carry
both bounds and use their maximum when composing certified operations.

For simplicity write all observations as one sequence indexed by $j$.
A tuple index $(i,k)$ can be flattened using a fixed numerical coding, or
one can keep rectangular tables with the same argument. Let
$\operatorname{Tab}(a,h)=[a(0),\ldots,a(h)]$. The zero extension of a
finite table gives a total input to any numerical program; correctness
conditions guarantee that no relevant access uses the extension.

\begin{definition}[Finite local contract]
\label{def:contract}
A stage $s$ of a construction is specified by:
\begin{itemize}
\item a primitive recursive numerical table program $F(s,p,X)$;
\item an evaluation bound $u(s,p)$ such that observation $p$ uses only the
      input entries indexed at most $u(s,p)$;
\item a Boolean finite property $C_s(h,X)$, using only the prefix through
      $h$ and downward closed under restriction to shorter prefixes;
\item a primitive recursive verification bound $v(s,h)$ and a proved
      finite implication
\begin{equation}
 C_s(v(s,h),X)=1
 \quad\Longrightarrow\quad
 C_{s+1}\bigl(h,T_s^h(X)\bigr)=1.             \label{eq:contract}
\end{equation}
\end{itemize}
Here $T_s^h(X)$ is the table of outputs $0\le p\le h$, evaluated from the
appropriate prefixes of $X$. The bound $v$ dominates the evaluation
bounds for those outputs. The predicates $C_s$ are values of one fixed
primitive recursive predicate $C(s,h,X)$. There is one \iPRA{}
derivation of \eqref{eq:contract} with $s,h$ and the metadata as free
numerical parameters, ordinarily obtained by Theorem~\ref{thm:safety}.
A separate external proof for every fixed $s$ does not meet this
uniformity requirement.
\end{definition}

Metadata can vary with the stage. It is included among the numerical
parameters and either computed exactly in advance or bounded by a fixed
primitive recursive recurrence. Metadata affecting the outputs is fixed
canonically from the original inputs and stage, independently of the
requested horizon, or is covered by the same prefix-locality condition.
In the bounded-metadata case replace the demands
by their bounded maxima over all metadata codes below the supplied cap;
substituting a cap into a nonmonotone demand is not sufficient.
If a particular interface requires extra output facts,
include their finite tests in $C_s$. We do not claim that every mathematical
property automatically has such a finite contract.

\begin{theorem}[Compilation of variable-depth observation constructions]
\label{thm:demand}
Suppose a fixed family of stages has the data of
Definition~\ref{def:contract}, uniformly primitive recursive in its
numerical stage and metadata. For any numerical number $N$ of stages
and desired final observation horizon $k$, there is one primitive
recursive table computation producing all final observations through
$k$ from a bounded prefix of the initial presentation. If that prefix
satisfies its initial finite contract, the output satisfies its final
finite contract. In applications a uniform derivation from the original
assumptions verifies the initial contract. The statement is uniform
in $N$ and $k$.

In particular, repeated application of \emph{one fixed} certified
operation does not require primitive recursion at function type, or a
universal evaluator for all primitive recursive programs.
\end{theorem}
\begin{proof}
With the starting metadata and $N$ fixed as numerical parameters, define
\[
 U_N(h)=\max\bigl( h,\{u(s,p):s<N,\ p\le h\},
                         \{v(s,p):s<N,\ p\le h\}\bigr).
\]
For $N=0$ the empty maxima contribute zero. This is a primitive recursive
function, nondecreasing in $h$, with $U_N(h)\ge h$. Stage-dependent metadata
is first computed by its fixed numerical bound recurrence and included in
the corresponding values of $u,v$.

Compute $d_0=k$ and $d_{r+1}=U_N(d_r)$ for $r<N$. For stage $s$ put
\begin{equation}
 h_s=d_{N-s}\qquad(0\le s\le N).               \label{eq:budgets}
\end{equation}
Thus $h_N=k$ and $h_s=U_N(h_{s+1})$. Initially tabulate the fixed input
observation function through $h_0$. Inductively form a new table
\begin{equation}
 X_{s+1}[p]=F\bigl(s,p,
          \operatorname{restrict}(X_s,u(s,p))\bigr),
 \qquad 0\le p\le h_{s+1}.                   \label{eq:table-build}
\end{equation}
The complete table is built by a bounded loop; the outer recursion is on
$s$ with $N,k$ and the already computed list of budgets fixed. Its state
is a \emph{number coding a finite table}. Every requested entry exists
since $u(s,p)\le h_s$.

All these operations are first-order primitive recursive. Nested
precision requests have become a forward computation of numerical
budgets followed by a forward computation of numerical tables. Even if
$U_N$ produces enormous exponent towers, its $N$-fold numerical iterate
is a primitive recursive function of $N,k$ and the starting metadata.
The number of stored observations is at most
$(N+1)(d_N+1)$. We do not claim a useful polynomial running time.

For the contract proof, use as the bounded invariant
\[
 s\le N\ \longrightarrow\ C_s(h_s,X_s)=1.
\]
The initial contract is a finite consequence of the input hypotheses.
Restriction supplies $C_s(v(s,h_{s+1}),\cdot)=1$ and
\eqref{eq:contract} gives the next row. Both the row construction and the
invariant are numerical. Theorem~\ref{thm:safety} applies, so there is
no induction on an unbounded statement about an entire output name.
\end{proof}

\paragraph{Coherence across requested horizons.}
Computing with two different final horizons must not change an output
already computed. This follows from deterministic prefix locality: if
$X,X'$ agree on every entry through $u(s,p)$, then their computed output
at $p$ is identical. An induction on the finite stage count shows
agreement of two table runs on their common relevant rows. Consequently
\[
 A_N(p):=X_N^{(N,p)}[p]
\]
is a single primitive recursive observation function. When a proof needs
several of its values simultaneously, compute one table at the maximum
verification horizon; locality identifies its entries with the individual
definition. This last step is essential: unrelated finite output tables
would not define a single observation procedure. Prefix consistency by
itself is not Cauchy regularity. When a regular name is required, its
pairwise finite regularity inequalities must be included in $C_N$;
arbitrary-horizon introduction then yields the full regularity schema.

The checked Lean theorem \texttt{compiled\_correct} verifies precisely
this comparison between the finite table implementation and the nested
semantic construction. Its companion \texttt{ofBody\_local} derives
locality from a numerical prefix-table body. The semantic construction
uses function-valued recursion only as an ambient Lean specification;
\eqref{eq:table-build} is the first-order arithmetic implementation.

\begin{corollary}[Eliminating function-valued iteration in a certified fragment]
\label{cor:first-order}
A construction written informally as
$A_{s+1}=T_s(A_s)$ admits a uniform first-order primitive recursive
implementation whenever its stage programs, metadata bounds, evaluation
demands, and finite verification contracts satisfy the theorem. Exact
algebraic operations, finite probability transformations, formal-series
coefficient operations, and precision-indexed analytic constructions can
share this same compiler.
\end{corollary}

\subsection{Finite error transport through a chain}

Many operations are verified only up to an arbitrarily small rational
error. A single error estimate can prevent induction on exact equality
of represented objects.

\begin{lemma}[Telescoping defects]
\label{lem:defects}
In a normed ring of finite rational data with a submultiplicative norm,
let $P_0,\ldots,P_N$ and $A_0,\ldots,A_{N-1}$ satisfy
\[
 \|P_s-A_sP_{s+1}\|\le\varepsilon_s,
 \qquad \|A_s\|\le G_s.
\]
Then
\begin{equation}
 \left\|P_0-\left(\prod_{s<N}A_s\right)P_N\right\|
 \le \sum_{s<N}\varepsilon_s\prod_{j<s}G_j.     \label{eq:defectsum}
\end{equation}
If every $\varepsilon_s\le\varepsilon$ and $G_s\le G$ with $G\ge1$,
the right side is at most $NG^N\varepsilon$.
\end{lemma}
\begin{proof}
For each $r\le N$, let $B_r=A_0\cdots A_{r-1}$, with $B_0=1$.
Subtract consecutive terms in $B_rP_r$:
\[
 B_rP_r-B_{r+1}P_{r+1}=B_r(P_r-A_rP_{r+1}).
\]
Summing this exact finite identity gives the difference on the left of
\eqref{eq:defectsum}. Repeated triangle inequality and
submultiplicativity give the asserted bound. The proof is induction on
finite coefficient lists, finite products, and finite sums; at no point
is the induction formula an equality of infinite names. For polynomials
use the coefficient $\ell^1$ norm, whose submultiplicativity is just the
finite convolution triangle inequality. If complex magnitudes are used
in explanatory notation, rational upper bounds or squared-norm
certificates give the finite numerical formulation.
\end{proof}

To prove an exact represented identity, first fix an output precision
$p$, then select a local tolerance satisfying
$NG^N\varepsilon\le2^{-p}$ and compile only the observations required at
that tolerance. Universal introduction in $p$, or
Lemma~\ref{lem:vanishing-error}, supplies the exact relation. This order
of operations also works for matrix products, compositions of approximate
isometries, telescoping integrals, and finite sums of probabilistic errors.

\subsection{Why the conditions are a real restriction}

There cannot be a total primitive recursive function $E(e,n)$ which
evaluates every unary primitive recursive function coded by $e$: if there
were, $d(n)=E(n,n)+1$ would itself be primitive recursive, with a code
$e_d$, and $d(e_d)=E(e_d,e_d)+1=d(e_d)+1$ would be impossible.
This is an external diagonal argument, not an assertion that \iPRA{}
proves its own soundness.

A uniformly bounded machine-step interpreter is different. Given a
program, its input, and an explicit fuel bound, one can perform that many
primitive recursive steps and return either an answer or a timeout. For
an individual fixed primitive recursive program, its numerical resource
bounds can be proved by its particular finite construction. Neither fact
creates a primitive recursive universal total evaluator.

Likewise, a merely pointwise statement that every stage eventually stops
does not satisfy Theorem~\ref{thm:demand}. The crucial input is the
\emph{uniform primitive recursive} numerical demand recurrence and its
proof. There is also no automatic conversion of an arbitrary Lean
correctness proof into a \texttt{PRAAx} proof. The finite safety theorem
requires actual base and step derivations in the object theory.

\subsection{Relation to proof interpretation}

The framework follows the established idea of extracting quantitative
content and majorants, but specializes it to explicit first-order
numerical compilation. Functional interpretation and majorization have a
much broader literature \cite{kohlenbach-notes,kohlenbach-point}. Their
higher-type settings must not be imported silently. In particular,
Kohlenbach's analysis of uniform weak K\"onig's lemma shows that apparently
small changes involving uniform selection and extensionality can greatly
increase proof-theoretic strength \cite{kohlenbach-uwkl}. Our applications
below exhibit the finite contracts directly and do not appeal to an
unspecified conservation theorem.

\section{An application: factorization uniformly in the degree}
\label{sec:fta-extension}

The repository's one-root FTA proof is uniform in the numerical degree,
but its complete-factorization paragraph deliberately claims only a
fixed-degree schema.  It asks for an additional bound on nested precision
requests before treating a variable number of deflations.  We supply that
bound here.  The input coefficient evaluator remains one fixed primitive
recursive program.  The result does not evaluate arbitrary codes of
primitive recursive programs.

The only substantive algebraic dependency is the repository's finite
rational descent lemma, whose proof uses bounded radius selection, five
fixed-length square-root bisections, and a finite binomial polynomial
\cite{repo-fta}.  We state its complete contract, reproduce the monic
one-root argument, and prove all additional estimates needed for uniform
factorization.  This is a human mathematical extension of the repository
proof, not an assertion that its outstanding deeply embedded Lean
derivation has already been constructed.

\subsection{Norms, finite names, and rounding}

For a Gaussian rational $z=x+iy$, put $|z|_1=|x|+|y|$.
Then $|zw|_1\le |z|_1|w|_1$, $|z|\le |z|_1\le 2|z|$, where
$|z|$ is the Euclidean norm.  Euclidean inequalities in the descent
argument are expressed arithmetically by squared rational comparisons.
For a finite polynomial $F$, put
\[
 \|F\|_1=\sum_j |[X^j]F|_1.
\]
Finite convolution and the triangle inequality prove
$\|FG\|_1\le\|F\|_1\|G\|_1$.

Let $\operatorname{rd}_v(z)$ round each real coordinate down to the
nearest multiple of $2^{-v}$.  This is a total primitive recursive
Gaussian-rational operation, and
\[
 |\operatorname{rd}_v(z)-z|_1\le 2^{1-v}.
 \tag{F1}\label{eq:fta-round}
\]
The convention for negative coordinates is ordinary rational floor.
Its error estimate follows directly from the two division-with-remainder
inequalities.  In particular $\operatorname{rd}_v(1)=1$.

For $n\ge1$ and $M\ge2$, call a coefficient table
$a_i(s)$, $0\le i\le n$, $0\le s\le h$, good through $h$ if
\begin{align*}
 a_n(s)&=1, & |a_i(s)|_1&\le M,\\
 |a_i(s)-a_i(t)|_1&\le 2^{-s}+2^{-t}
 &&(i\le n,\ s,t\le h).
\end{align*}
All quantifiers in this definition are bounded.  Its characteristic
function is a finite Boolean conjunction, so goodness is expressed by
one primitive recursive equation.  An infinite name is merely convenient
notation for the corresponding family of finite tables.

The use of $|\cdot|_1$ and an explicit coefficient cap is harmless for
the repository's fast Euclidean names.  Given such names $a_i$, replace
the lower coefficients by
\[
 \widetilde a_i(s)=\operatorname{rd}_{s+4}(a_i(s+3)),
 \qquad \widetilde a_n(s)=1.
\]
The new names represent the same coefficients and satisfy the displayed
$|\cdot|_1$ regularity.  Indeed the two source-name errors contribute at
most $\tfrac14(2^{-s}+2^{-t})$, and rounding contributes at most
$\tfrac18(2^{-s}+2^{-t})$.  A suitable cap is
\[
 M=4+\left\lceil\max_{i<n}|a_i(0)|_1\right\rceil.
\]
Here monicity is supplied as an exact presentation of the top coefficient;
we do not try to decide whether an arbitrary coefficient name equals $1$.

\subsection{A uniform finite root routine with an explicit query bound}

The following finite rational lemma is the imported result.

\begin{lemma}[Finite rational descent contract]\label{lem:fta-descent-contract}
For $n\ge1$, put $D=2^n4^{2n^2}$ and $q=1-(4D)^{-1}$.
There is one primitive recursive routine
$\operatorname{Descent}(n,P,z)$ such that for every nonconstant Gaussian
rational polynomial $P$ of degree at most $n$, it returns Gaussian rational
$\Delta$ and nonnegative rational $r_0$ with
\[
 |P(z+\Delta)|\le q|P(z)|,\qquad |\Delta|\le r_0,
 \qquad |b_i|r_0^i\le |P(z)|\quad(1\le i\le n),
\]
where $P(z+W)=\sum_i b_iW^i$.
All loops and searches in this routine have explicit primitive recursive
bounds in the input codes.  Failed searches and malformed inputs have
fixed default outputs.
\end{lemma}

For clarity, the bound is uniform in $n$; it is not an external family
of routines indexed by fixed $n$.  The source proves this lemma by
testing a finite range of dyadic radii, finding a repeated maximizing
index among $2n+1$ scales, and constructing the appropriate direction by
the finite binomial identity.  These operations introduce no new calls
to the coefficient-name evaluator.

Fix $n,M$ as above and define
\begin{align*}
 B&=M+2, &D&=2^n4^{2n^2}, &K&=8nD,\\
 \rho&=1-K^{-1}, &Q&=1-(8D)^{-1}, &R&=BK+1,\\
 L&=2+\operatorname{bl}(16D(n+1)R^n),
 &H&=\operatorname{bl}(BK),\\
 J(t)&=K(t+H+2), &U(n,M,t)&=L+J(t).
 \tag{F2}\label{eq:fta-constants}
\end{align*}
Here $\operatorname{bl}(b)$ is the binary digit count, so
$b<2^{\operatorname{bl}(b)}$ for $b>0$.
The symbols $B,D,K,R,L,H,J$ in subsequent formulas refer to these explicit
functions of $n,M$.

For the given coefficient table set
\[
 P_j(X)=\sum_{i=0}^n a_i(L+j)X^i,
 \qquad z_0=0,
 \qquad z_{j+1}=z_j+\operatorname{Descent}(n,P_j,z_j).\Delta.
 \tag{F3}\label{eq:fta-trajectory}
\]
This is one fixed primitive recursion.  Define
\[
 \alpha(t)=z_{J(t)},\qquad
 r(t)=\operatorname{rd}_{t+4}(\alpha(t+2)),\qquad A=2R+1.
 \tag{F4}\label{eq:fta-root-normalized}
\]

\begin{lemma}[Finite monic root routine]\label{lem:fta-root-table}
For an output horizon $h$, goodness of the input table through
$U(n,M,h+2)$ suffices to prove
\begin{align*}
 |r(t)|_1&\le A &&(t\le h),\\
 |r(t)-r(u)|_1&\le 2^{-t}+2^{-u} &&(t,u\le h).
\end{align*}
The value $r(t)$ reads coefficient entries only at indices at most
$U(n,M,t+2)$, and it is unchanged if the input table is extended.
In addition, setting
\[
 E=4(n+1)^2(M+1)(A+1)^{n+1},
 \tag{F5}\label{eq:fta-E}
\]
the same finite hypotheses give
\[
 \left|\sum_{i=0}^n a_i(t)r(t)^i\right|_1\le E2^{-t}
 \qquad(t\le h).
 \tag{F6}\label{eq:fta-same-precision-residual}
\]
Every assertion has a quantifier-free-induction proof after the displayed
finite horizons are fixed.
\end{lemma}

\begin{proof}
The leading coefficient of every $P_j$ is $1$.  The radius certificate
in Lemma~\ref{lem:fta-descent-contract} therefore gives
$r_0^n\le |P_j(z_j)|$.  Finite Bernoulli arithmetic gives
$Q\le\rho^n$, and hence
\[
 |P_j(z_j)|\le BQ^j\quad\Longrightarrow\quad
 |\Delta_j|\le B\rho^j.
\]
Indeed $(B\rho^j)^n\ge B\rho^{nj}\ge BQ^j$, and positive integer
powers reflect order on nonnegative rationals.  This argument avoids
computing any $n$th root.

Induct simultaneously on the finite inequalities
\[
 |P_j(z_j)|\le BQ^j,
 \qquad |z_j|\le B\sum_{v<j}\rho^v\le BK<R.
 \tag{F7}\label{eq:fta-monic-invariant}
\]
The base case uses $|a_0(L)|\le M\le B$.  The step bound just proved
establishes the position bound at $j+1$.  At any point $|x|\le R$, the
adjacent coefficient-name comparisons give
\[
 |P_{j+1}(x)-P_j(x)|
 \le\frac32(n+1)R^n2^{-L-j}
 \le\frac{3}{32D}2^{-j}
 \le\frac{1}{8D}BQ^j.
\]
The choice of $L$ is stronger than the inequality used here.  Rational
descent then changes the residual bound by the factor
$q+(8D)^{-1}=Q$.  This proves \eqref{eq:fta-monic-invariant}.

Finite geometric summation gives
\[
 |z_v-z_u|\le BK\rho^u\qquad(v\ge u).
\]
Also $\rho^K\le\tfrac12$: apply finite Bernoulli to
$(1+(K-1)^{-1})^K$.  Thus
\[
 |\alpha(v)-\alpha(u)|\le 2^{-u-2}\qquad(v\ge u).
 \tag{F8}\label{eq:fta-raw-root-cauchy}
\]
If $s\ge U(n,M,t)$ and the input table also contains row $s$, comparison
with $P_{J(t)}$ gives
\[
 \left|\sum_{i=0}^n a_i(s)\alpha(t)^i\right|
 \le BQ^{J(t)}+2(n+1)R^n2^{-L-J(t)}
 \le 2^{-t}.
 \tag{F9}\label{eq:fta-raw-root-residual}
\]
This is a finite statement about the specified $s$, not a limit argument.

Rounding in \eqref{eq:fta-root-normalized} adds $|\cdot|_1$ error at
most $2^{-t-3}$.  Combining this with twice the Euclidean bound
\eqref{eq:fta-raw-root-cauchy} proves regularity of $r$ (in fact with
slack), and $|r(t)|_1\le 2R+1=A$.

For \eqref{eq:fta-same-precision-residual}, use
\eqref{eq:fta-raw-root-residual} at $t+2$ and
$s=U(n,M,t+2)$.  Replacing row $s$ by row $t$ costs at most
$2(n+1)A^n2^{-t}$ in $|\cdot|_1$.  The residual in
\eqref{eq:fta-raw-root-residual} costs at most $2^{-t-1}$ in that norm.
Finally replacing $\alpha(t+2)$ by $r(t)$ costs at most
\[
 \frac18 M\left(\sum_{i=1}^n iA^{i-1}\right)2^{-t}
\]
by the finite difference-of-powers identity.  Their sum is less than
$E2^{-t}$ for \eqref{eq:fta-E}.

The largest coefficient index used in these arguments is the displayed
$U(n,M,h+2)$.  In the trajectory, the invariant is a pair of rational
inequalities; its Euclidean components can be squared.  Geometric-sum
and coefficient comparisons are finite folds.  A Boolean conjunction
stores all required assertions through the fixed horizon.  Quantifier-free
induction therefore proves the finite contract.  Prefix compatibility
follows because each run follows exactly the recursion
\eqref{eq:fta-trajectory}, with no decision based on table length.
\end{proof}

\subsection{A normalized deflator and its complete finite contract}

Use the following additional integer bounds:
\begin{align*}
 M^+&=2+(n+1)(M+1)(A+1)^n,\\
 a&=\operatorname{bl}(4E),\\
 D^+&=4(n+1)(M^++1)(A+1),\\
 V(n,M,h)&=U(n,M,h+a+2).
 \tag{F10}\label{eq:fta-deflation-budgets}
\end{align*}
The letter $a$ in this display denotes the precision shift, not a
coefficient.  For $0\le j<n$ put
\begin{align*}
 q_j(v)&=\sum_{\ell=j+1}^n a_\ell(v)r(v)^{\ell-j-1},\\
 b_j(t)&=\operatorname{rd}_{t+4}(q_j(t+a)),\\
 P(t;X)&=\sum_{j=0}^n a_j(t)X^j,
 &P^+(t;X)&=\sum_{j=0}^{n-1} b_j(t)X^j.
 \tag{F11}\label{eq:fta-deflator}
\end{align*}
These are finite sums, and $b_{n-1}(t)=1$ exactly.

\begin{lemma}[Finite deflation contract]\label{lem:fta-deflation-table}
If the input coefficient table is good through $V(n,M,h)$, then the
output table $b_j(t)$ is monic, bounded by $M^+$ in $|\cdot|_1$, and
regular through $h$.  The normalized root is regular and bounded by $A$
through $h$.  For every $t\le h$,
\[
 \|P(t;X)-(X-r(t))P^+(t;X)\|_1\le D^+2^{-t}.
 \tag{F12}\label{eq:fta-defect}
\]
The complete calculation and every finite verification use only input
indices at most $V(n,M,h)$.  Extending the input table changes none of the
already computed output entries.
\end{lemma}

\begin{proof}
Each $q_j$ has at most $n$ terms.  Hence
\[
 |q_j(v)|_1\le nM(A+1)^n<M^+-1.
\]
Rounding contributes at most $2^{-t-3}\le1$, proving the cap.
For any two available indices $v,w$, subtract corresponding terms and
apply the difference-of-powers identity.  Input and root regularity give
\[
 |q_j(v)-q_j(w)|_1
 \le\left(\sum_{r<n} A^r+M\sum_{r<n}rA^{r-1}\right)
       (2^{-v}+2^{-w})
 \le E(2^{-v}+2^{-w}),
\]
where the $r=0$ term of the second sum is $0$.
Since $E2^{-a}<1/4$, after the index shift and rounding this is at most
$\tfrac38(2^{-t}+2^{-u})$ for $b_j(t),b_j(u)$, which implies the
required regularity.

Set $v=t+a$ and $Q_v(X)=\sum_{j<n}q_j(v)X^j$.  Synthetic division is
the exact finite identity
\[
 P(v;X)=(X-r(v))Q_v(X)+P(v;r(v)).
\]
It follows by cancelling adjacent coefficients, including the constant
coefficient.  The defect in \eqref{eq:fta-defect} is the sum of four
terms:
\begin{align*}
 &[P(t;X)-P(v;X)] + P(v;r(v))\\
 &\quad +(X-r(v))[Q_v(X)-P^+(t;X)]
       +(r(t)-r(v))P^+(t;X).
\end{align*}
Their coefficient norms are bounded respectively by
\[
 2(n+1)2^{-t},\qquad
 E2^{-v}<\tfrac14 2^{-t},\qquad
 (A+1)n2^{-t-3},\qquad
 2nM^+2^{-t}.
\]
Their sum is at most $D^+2^{-t}$.

The residual and root regularity at indices at most $h+a$ require the
source root contract through $U(n,M,h+a+2)=V(n,M,h)$.  All other
coefficient queries are smaller.  Thus $V$ is a verification bound as
well as an evaluation bound.  The proof consists of finite coefficient
identities, bounded sums, and rational comparisons; no induction on a
statement about all precisions has been used.
\end{proof}

\subsection{Variable-degree iteration by finite tables}

Let the input degree be $d\ge1$ and initial cap be $M_0=M$.
Compute the numeric caps and constants by the single primitive recursion
\[
 n_i=d-i,\qquad
 M_{i+1}=M^+(n_i,M_i)\quad(0\le i<d).
 \tag{F13}\label{eq:fta-cap-recursion}
\]
Write $A_i,D_i^+,V_i$ for the corresponding functions in
\eqref{eq:fta-deflation-budgets}.  No coefficient evaluation is needed
to compute these bounds.

For a desired output horizon $t$, determine input horizons backwards:
\[
 h_d=t,\qquad h_i=V_i(h_{i+1})\quad(i=d-1,\ldots,0).
 \tag{F14}\label{eq:fta-backwards}
\]
This is a primitive recursion over a reversed finite index list.  It is
not recursion on a program code.  Each $V_i$ is monotone in its last
argument and is at least that argument.

Generate the original coefficient table through $h_0$ by the one fixed
PR input evaluator.  At stage $i$, run the finite deflator on its current
table to obtain its root entries and next coefficient entries through
$h_{i+1}$.  After $d$ stages the residual polynomial is exactly $1$.
Each stage is one application of the same fixed PR table routine to
numeric inputs $n_i,M_i,h_{i+1}$ and the current finite table.  Therefore
the entire staged calculation is primitive recursive uniformly in $d$.

For an output root at index $i<d$ and precision $p$, use this algorithm
with terminal horizon $p$ and select the $i$th root entry at $p$.
Call it $r_i(p)$.  Prefix compatibility in the two local lemmas implies
that computing with any larger terminal horizon gives the identical
rational entry.  To verify this assertion, fix the two finite runs and
induct on their common stage index: every lookup needed for a common
output row has the same input value by the induction hypothesis, so the
deterministic recursion returns the same value.  This finite argument establishes prefix consistency. For two output
precisions $p,q$, run the construction at terminal horizon $\max(p,q)$,
apply root regularity there, and identify its entries with the separately
defined outputs. Thus the outputs are regular names; no compatible-root
choice principle is needed.

\subsection{The uniform factorization theorem}

Set
\[
 G=1+\max_{i<d}A_i,\qquad
 W=1+\max_{i<d}D_i^+,\qquad
 c=\operatorname{bl}(1+dWG^d).
 \tag{F15}\label{eq:fta-final-budget}
\]

\begin{theorem}[Degree-uniform constructive factorization schema]
\label{thm:fta-uniform-factorization}
Fix one primitive recursive coefficient evaluator
$a(d,\mathbf p,i,s)$, with natural parameters $d,\mathbf p,i,s$.
There is one primitive recursive root evaluator
$r(d,M,\mathbf p,i,s)$, and the explicit primitive recursive bound
$c=c(d,M)$ above, such that iPRA proves the following.
If the input is monic of degree $d\ge1$, its coefficient names are
regular, and $|a(d,\mathbf p,i,s)|_1\le M$ for $i\le d$, then the
constructed $r_i$ are regular coefficient names and, for every $k$,
\[
 \left\|
   \sum_{j=0}^d a(d,\mathbf p,j,k+c)X^j
       -\prod_{i<d}(X-r_i(k+c))
 \right\|_1\le2^{-k}.
 \tag{F16}\label{eq:fta-uniform-factorization}
\]
In particular these names represent a full factorization, with
multiplicities.  All degrees are handled by the same numeric program,
not by a separately compiled evaluator for each degree.
\end{theorem}

\begin{proof}
Fix $k$ and put $t=k+c$.  Use the finite schedule
\eqref{eq:fta-backwards}.  The global input assumptions establish the
finite good-table condition at $h_0$ by bounded Boolean-fold induction.
Lemma~\ref{lem:fta-deflation-table} propagates the good-table conditions
through the $d$ stages.  Keep only the finite coefficients at the common
precision $t$.  Write $P_i$ for the degree-$d-i$ sampled polynomial,
$L_i=X-r_i(t)$, and $e_i=P_i-L_iP_{i+1}$.
Then $P_d=1$, $\|L_i\|_1\le G$, and
$\|e_i\|_1\le W2^{-t}$.  Repeated substitution gives the exact identity
\[
 P_0-\prod_{i<d}L_i
   =\sum_{i<d}\left(\prod_{j<i}L_j\right)e_i.
 \tag{F17}\label{eq:fta-telescope}
\]
Finite convolution bounds therefore give
\[
 \left\|P_0-\prod_{i<d}L_i\right\|_1
 \le W2^{-t}\sum_{i<d}G^i
 \le dWG^d2^{-t}\le2^{-k}.
\]
This proves \eqref{eq:fta-uniform-factorization}.

For the logical-strength claim, the stage invariant comprises the Boolean
good-table checks through the fixed $h_i$, the finitely many local
defect inequalities, and the finite coefficient identity
\eqref{eq:fta-telescope} through the current stage.  Every quantified
index is bounded by $d$, a listed $h_i$, or a coded array length.
Their Booleanization is quantifier-free.  Quantifier-free induction is
therefore enough.  Universal introduction over $k$ follows only after
this finite proof.  There is no induction on represented-complex equality
or on the assertion that a polynomial factors at all precisions.
\end{proof}

The dependence on the initial evaluator is exactly the dependence allowed
by the original FTA schema: once that PR definition is fixed, the table
generator is a composition with it.  For a finite rational-polynomial
code, constant coefficient names give a single first-order specialization
uniform in that code and its degree.  For an apart leading coefficient,
the usual certified reciprocal and a fixed precision shift reduce to the
monic formulation.  This does not decide which coefficient is leading
when the degree has not been supplied.

\subsection{Resource bounds and the remaining Lean work}

The query bound and verification bound have different roles.
The raw root at precision $t$ reads through $U(n,M,t)$; the normalized
root reads through $U(n,M,t+2)$.  A deflation through output horizon $h$
is verified from goodness through $V(n,M,h)$, which includes shifted
coefficient comparisons and the same-precision residual estimate.
Composing only the smaller query bound would not justify the theorem.

No non-primitive-recursive search is hidden in rational encoding sizes.
The fixed root table routine performs exactly $J(t)$ descent iterations;
each descent has bounded loops in its current finite codes, as spelled
out in the imported lemma.  The complete table map is therefore a fixed
PR function.  Iteration over $d$ takes table \emph{values} as state; it
never forms a growing PR program and asks a universal evaluator to run it.

If a numerical arithmetic-operation envelope is desired, first use the dyadic
normalization above.  A finite input table then has a primitive recursive
code bound $Z(n,M,h)$: with any fixed tuple encoding, bound the signed
dyadic numerators by $(M+1)2^{h+4}$ and the denominators by $2^{h+4}$,
then bound the $(n+1)(h+1)$ entries by a positional digit code.  Instrument
the one fixed table routine with an arithmetic-operation counter, using
the same compositions and recursions.  Its counter function $T$ is PR. This counts operations of the fixed table
transformer; the cost of generating the original coefficient table remains
the cost of its particular fixed evaluator. No bit-cost estimate is claimed.
The explicit bounded maximum
\[
 \max_{z\le Z(n,M,V(n,M,h))}
      T(n,M,h,z)
\]
is a PR envelope, with malformed tables totalized as above.  This maximum
is over finite input data for one fixed routine, not over arbitrary
programs.  An independent complexity claim stronger than primitive
recursiveness is not needed.

A Lean implementation should define the finite table data, $U,V$, the
cap recursion, rounding, and polynomial coefficient norm first.  The
critical new lemmas are the finite normalized-root contract,
the finite deflator contract, prefix compatibility, and the coefficient
telescoping bound.  Their precise statements are
Lemmas~\ref{lem:fta-root-table} and~\ref{lem:fta-deflation-table},
\eqref{eq:fta-telescope}, and
Theorem~\ref{thm:fta-uniform-factorization}.  Their arithmetic versions
must then construct object-language proofs using the Booleanized finite
invariants just described.  The repository's local descent derivation
remains an imported human proof until that object derivation is finished;
a Lean theorem that merely accepts the descent contract as a hypothesis
must be labelled conditional.

\section{Quantitative compactness without a compactness axiom}
\label{sec:quantitative-analysis}

This section supplies a reusable replacement for three recurring analytical
steps: taking an extremal value on a compact set, obtaining an actual point
from approximate solutions, and inferring local continuity from convexity.
The input data and the precision schedules are part of the theorems.  The
results do not infer a primitive-recursive modulus from arbitrary qualitative
continuity.  The distinction matters in the calculus and Taylor--Cauchy
packages, which explicitly supply their continuity and derivative contracts
\cite{repo-calculus,repo-taylor}.

\subsection{The presentation convention}

For definiteness work on
\[
 K_d=[0,1]^d\cap\Q^d,\qquad
 d_\infty(x,y)=\max_{i<d}|x_i-y_i|.
\]
For $d=0$ the domain contains one empty tuple and its distance is zero.
The same arguments apply to other rational domains when finite nets and
their rounding maps are supplied.  A scalar presentation consists of fixed
primitive-recursive rational functions
\[
 A(x,p),\qquad \mu(k),
\]
with the following two contracts, for rational points in $K_d$:
\begin{align}
 |A(x,p)-A(x,q)|&\le 2^{-p}+2^{-q},
       \label{eq:analysis-coherence}\\
 d_\infty(x,y)\le2^{-\mu(k)}
 &\ \Longrightarrow\ 
 |A(x,p)-A(y,q)|\le2^{-k}+2^{-p}+2^{-q}.
       \label{eq:analysis-continuity}
\end{align}
We write $f(x)$ for this coherent rational name.  Expressions involving
$f(x)$ are notation for rational inequalities at arbitrary requested
precisions, rather than a new sort of completed real numbers in \iPRA.
For example, $f(x)\le r$ means
$A(x,p)\le r+2^{-p}$ for every $p$.  For sums and products one first
allocates the finite error budget and then obtains an equivalent family of
rational inequalities.

Whenever we remove an arbitrarily small positive error from a rational
inequality, the operation has an elementary finite explanation.  If a
positive rational gap is $u/v$, with positive integers $u,v$, a dyadic error
smaller than $1/v$ contradicts that gap.  A suitable exponent is obtained by
bounded search below $v+1$.  Thus arguments such as ``for every precision,
therefore the limiting inequality holds'' can be translated into a
particular rational contradiction at a computable precision.  They are not
applications of a completeness principle.

The functions $A$ and $\mu$ are fixed syntax in the object theory, or are
specific functions already present in a definitional extension.  The
theorem is a schema over such presentations.  In particular, there is no
primitive-recursive universal evaluator for arbitrary primitive-recursive
program codes hidden in the phrase ``evaluate $A$''.

\subsection{A finite-net extremum theorem}

\begin{theorem}[A primitive-recursive infimum name]
\label{thm:finite-net-infimum}
From \eqref{eq:analysis-coherence}--\eqref{eq:analysis-continuity} one can
construct a primitive-recursive sequence of rational points $a_k\in K_d$
and nested rational intervals $[L^*_k,U^*_k]$ satisfying
\begin{align}
 0\le U^*_k-L^*_k&\le\tfrac34\,2^{-k},
       \label{eq:infimum-width}\\
 L^*_k&\le f(x)\quad(x\in K_d),
       \label{eq:infimum-lower}\\
 f(a_k)&\le V+\tfrac34\,2^{-k},
       \label{eq:approx-minimizer}
\end{align}
where $V$ is the name given by these nested intervals.  The name $V$ is a
lower bound for all $f(x)$ and is the unique named value with this lower
bound and approximate-attainment property.

The construction uses
\[
 n_k=\mu(k+2),\qquad
 N_k=(2^{n_k}+1)^d
\]
grid evaluations, all at precision $k+2$.  It does not assert that the
points $a_k$ form a Cauchy name.
\end{theorem}

\begin{proof}
Let $G_n$ be the full dyadic grid whose coordinates are
$0,2^{-n},\ldots,1$.  Enumerate it in a fixed lexicographic order.
For $x\in K_d$, rounding each coordinate down gives a grid point
$\pi_n(x)$ with
\[
 d_\infty(x,\pi_n(x))\le2^{-n}.
\]
Both the enumeration and rounding are rational, bounded computations.

Put $\varepsilon_k=2^{-k}$.  Evaluate every $g\in G_{n_k}$ at
precision $k+2$, let $m_k$ be the least sampled value, and let $a_k$ be
the first point attaining it.  Rational order makes the tie-breaking
decidable.  For every rational $x\in K_d$,
\begin{align*}
 m_k
 &\le A(\pi_{n_k}(x),k+2)\\
 &\le f(\pi_{n_k}(x))+\tfrac14\varepsilon_k\\
 &\le f(x)+\tfrac12\varepsilon_k.
\end{align*}
Also
$f(a_k)\le m_k+\varepsilon_k/4$.
Define
\[
 L_k=m_k-\tfrac12\varepsilon_k,
 \qquad U_k=m_k+\tfrac14\varepsilon_k.
\]
Consequently $L_k\le f(x)$ for every rational $x$, whereas
$f(a_k)\le U_k$.  In particular, for every $i,j$,
\[
 L_i\le f(a_j)\le U_j.
\]
The endpoints are rational, so removing the arbitrarily small evaluation
errors gives the literal rational inequality $L_i\le U_j$.

Now take finite running intersections:
\[
 L^*_k=\max_{i\le k}L_i,
 \qquad U^*_k=\min_{i\le k}U_i.
\]
The cross inequalities just proved show $L^*_k\le U^*_k$.  The
intervals are nested, and their width is at most
$U_k-L_k=3\varepsilon_k/4$.  This proves
\eqref{eq:infimum-width} and defines the name $V$ without choosing a
limit. Explicitly, a scalar sampler for this name is
$V(k)=(L^*_k+U^*_k)/2$. For $l\ge k$ both midpoints are in the
$k$th interval, so $|V(l)-V(k)|\le3\,2^{-k}/4\le2^{-k}+2^{-l}$.
Thus it is regular in the sense of \eqref{eq:RN}.
All its lower endpoints are lower bounds for $f$. In particular,
$V(q)\le f(x)+3\,2^{-q}/8$; after allowing the sampling error of
$f(x)$, this is the regular-name lower-bound inequality at index $q$.
Moreover,
as $V\ge L_k$ and $f(a_k)\le U_k$,
\[
 f(a_k)\le V+(U_k-L_k)
          =V+\tfrac34\varepsilon_k.
\]
These relations can each be unfolded at any finite requested precision.

Finally suppose $W$ is another lower-bound name with approximate
attainment.  Since $W\le f(a_k)\le V+3\varepsilon_k/4$ for every
$k$, we have $W\le V$ as a named inequality.  Apply the corresponding
approximate-attainment property of $W$ to obtain $V\le W$.  No
extremizing point or subsequence is used in this comparison.
\end{proof}

The maximum is obtained by applying the construction to $-f$.  A
rational box other than $[0,1]^d$ is handled by an affine change of
variables with its explicit distance factor.  More generally, only a
primitive-recursive finite net, a rounding map, and a modulus were used.
Thus the result is a general finite-net theorem rather than a special
property of cubes.

The construction resolves a useful part of an appeal to the extreme-value
theorem: a supremum or infimum \emph{value}, together with points as close
to the value as requested, has a direct finite realization.  It does not
resolve the separate task of returning a coherent extremizing point.
For example, choosing alternating points $0,1$ for the constant function
returns exact minimizers at every stage, but those choices are not a
Cauchy sequence.  A rule that identifies approximate solutions at nearby
precisions is necessary for the next step.

\subsection{Quantitative uniqueness makes approximate witnesses coherent}

\begin{theorem}[Coherent extraction from approximate minima]
\label{thm:coherent-minimizer}
In the setting of Theorem~\ref{thm:finite-net-infimum}, suppose a fixed
primitive-recursive function $u$ satisfies the quantitative uniqueness
contract
\begin{equation}
 \left.
 \begin{array}{c}
 f(x)\le V+2^{-u(k)}\\
 f(y)\le V+2^{-u(k)}
 \end{array}\right\}
 \quad\Longrightarrow\quad d_\infty(x,y)\le2^{-k}.
 \label{eq:quantitative-uniqueness}
\end{equation}
Set $v(k)=k+\max_{j\le k}u(j)$ and $x_k=a_{v(k)}$.
Then
\[
 d_\infty(x_k,x_l)\le2^{-\min(k,l)}
                    \le2^{-k}+2^{-l}.
\]
The $x_k$ therefore give a primitive-recursive point name $X$ in the
closed cube.  Its canonical evaluation satisfies $f(X)=V$, and it is
the unique point name with that property.
\end{theorem}

\begin{proof}
If $k\le l$, both $v(k)$ and $v(l)$ are at least $u(k)$.
Equation~\eqref{eq:approx-minimizer} shows that $x_k,x_l$ both meet
the hypotheses of \eqref{eq:quantitative-uniqueness} at precision $k$.
This proves the displayed Cauchy estimate.  All coordinates of all
approximants lie between zero and one, so the same closed inequalities
hold of the point name.

For completeness, a canonical evaluation procedure is explicit.  Replace
$\mu$ by its monotone majorant if necessary and put
\[
 r(p)=p+2+\max_{j\le p+2}\mu(j),\qquad
 B(p)=A(x_{r(p)},p+2).
\]
The distance bound between $x_{r(p)}$ and $x_{r(q)}$, followed by
\eqref{eq:analysis-continuity}, proves coherence of $B$ with ample
dyadic slack.  This defines $f(X)$.  Further,
$v(r(p))\ge r(p)\ge p+2$, so the optimum gap of $x_{r(p)}$ is at
most $2^{-p-2}$.  The evaluation error is also $2^{-p-2}$.
Hence $B(p)$ differs from $V$ by at most $2^{-p-1}$, interpreted using
the interval name for $V$.  This proves equality.

If another point name $Y$ has value $V$, choose rational approximants
to $X,Y$ fine enough for the continuity modulus and for the tolerance in
\eqref{eq:quantitative-uniqueness}.  The equality of the values supplies
the required approximate-minimum premises.  The uniqueness contract and
the coordinate errors show $d_\infty(X,Y)$ is below any prescribed
dyadic tolerance.  Thus $X=Y$ as point names.
\end{proof}

\begin{corollary}[A complete strongly convex minimization scheme]
\label{cor:strong-convexity}
Suppose a positive rational $\lambda$ is supplied and
\begin{equation}
 f\!\left(\frac{x+y}{2}\right)
 \le\frac{f(x)+f(y)}2
       -\frac{\lambda}{8}\,d_\infty(x,y)^2
 \quad(x,y\in K_d).
 \label{eq:midpoint-drop}
\end{equation}
Let $c$ be a natural number with $\lambda2^c\ge8$.
Then $u(k)=2k+c$ satisfies \eqref{eq:quantitative-uniqueness}.
In particular the simpler sequence
\[
 x_k=a_{2k+c}
\]
already gives a coherent name of the unique minimizer.
\end{corollary}

\begin{proof}
The midpoint belongs to $K_d$, so the lower-bound property of $V$
and \eqref{eq:midpoint-drop} give
\[
 V\le f\!\left(\frac{x+y}{2}\right)
   \le V+2^{-u(k)}
                 -\frac{\lambda}{8}d_\infty(x,y)^2.
\]
Thus
\[
 d_\infty(x,y)^2
 \le\frac8\lambda\,2^{-2k-c}\le2^{-2k}.
\]
The distance is rational and nonnegative.  Rational order therefore
gives $d_\infty(x,y)\le2^{-k}$; a positive contrary gap would remain
positive after squaring.  This is precisely the uniqueness contract.
The schedule $2k+c$ is already monotone, and the approximate minima
on this schedule have a vanishing value gap, so the preceding proof
applies directly.

The integer $c$ need not be found by unbounded logarithmic search.
Write $\lambda=a/b$ with positive integers $a,b$.  Searching
$0\le c\le8b$ suffices, since $2^{8b}\ge8b\ge8b/a$.
Consequently every step of the stated schedule is primitive recursive.
\end{proof}

The usual Euclidean $\lambda$-strong-convexity inequality implies
\eqref{eq:midpoint-drop}, because
$\|x-y\|_2^2\ge d_\infty(x,y)^2$.  One may also prove
\eqref{eq:midpoint-drop} directly for a particular rationally presented
objective, without introducing derivatives or Hessians.  For example,
$f(x)=\sum_{i<d}(x_i-b_i)^2$ has $\lambda=2$, by expansion of squares.
The theorem then supplies a complete construction of the minimizing
point, including boundary cases, from finite rational minimization.
This is stronger than merely replacing an extremum by a finite
certificate: it constructs the coherent point and proves its optimality.

For equations one can use the same pattern with a nonnegative residual
$\rho$.  A finite algorithm producing $\rho(a_k)\le2^{-k}$ supplies
approximate witnesses.  A quantitative isolation contract
\[
 \rho(x),\rho(y)\le2^{-u(k)}
 \ \Longrightarrow\ d(x,y)\le2^{-k}
\]
supplies coherence.  Neither ingredient follows from the other.  In
particular, multiple roots may violate this isolation contract.  The
existing FTA package instead controls the movement of its descent
sequence and already distinguishes that argument from unrelated
approximate zeros \cite{repo-fta}.  Our theorem does not replace its
multiple-root construction by an unjustified uniqueness assumption.

\subsection{Removing the extra candidate modulus in Bohr--Mollerup}
\label{sec:bm-modulus-repair}

The Bohr--Mollerup report asks whether the primitive-recursive local
modulus of the candidate can be derived rather than supplied
\cite{repo-bm}.  For its rationally presented functions, the answer is
yes.  In fact, recurrence is unnecessary for this part.

\begin{lemma}[Positive multiplicative convexity implies ordinary convexity]
\label{lem:positive-convexity}
Let $F$ assign strongly positive regular names to positive rationals.
Here strong positivity includes a fixed primitive recursive witness
function giving a positive rational lower bound at each rational
argument, as in the repository's presentation; it is not merely an
unwitnessed assertion of pointwise positivity.
Suppose that, for $s\ge1$, $0\le j\le s$, and
$z=((s-j)x+jy)/s$,
\begin{equation}
 F(z)^s\le F(x)^{s-j}F(y)^j.
 \label{eq:rlc-framework}
\end{equation}
Then
\[
 F(z)\le\frac{s-j}{s}F(x)+\frac{j}{s}F(y).
\]
The implication has a finite arithmetic proof with
primitive-recursive precision demands.
\end{lemma}

\begin{proof}
Finite weighted arithmetic--geometric mean gives
\[
 F(x)^{s-j}F(y)^j
 \le\left(\frac{(s-j)F(x)+jF(y)}s\right)^s.
\]
This is the finite AM--GM theorem already developed in the same
Bohr--Mollerup package, applied to a list with $s-j$ copies of $F(x)$
and $j$ copies of $F(y)$.  Its elementary argument can be recalled
briefly.  For two positive entries, $(u-v)^2\ge0$ proves the result.
For a list of power-of-two length, repeatedly split it into two equal
blocks.  Applying the same inequality to the means of the two blocks,
and raising it to their common length, proves the next level of the
binary product tree.  For a list of arbitrary length $s$, let $m$ be
the least power of two at least $s$, and let $a$ be the mean of the
original list.  Pad the list by $m-s$ copies of $a$.  Its mean is still
$a$, so the power-of-two case gives
$\bigl(\prod_{i<s}u_i\bigr)a^{m-s}\le a^m$.
Cancellation of the positive factor $a^{m-s}$ proves the required
$s$-entry inequality.  There are at most $2s$ leaves in the binary
product tree.  This padding construction avoids a separate downward
induction through all intervening lengths.  At every node the
invariant is an inequality of finite products and rational means;
its finite arithmetic version is proved by a Boolean prefix invariant
over the tree in postorder.

Combine AM--GM with \eqref{eq:rlc-framework} and cancel the common
positive power.  This last step does not decide the order of two real
names.  To verify the desired order at tolerance $\varepsilon>0$,
a putative contrary rational gap of that size forces a positive gap
of at least $\varepsilon^s$ between the $s$th powers, by the finite
factorization of $u^s-v^s$ or the binomial theorem.  Requesting the
power evaluations to error less than half that gap gives a rational
contradiction.  The magnitude bounds needed to evaluate the powers
come from coarse input intervals; $s$, $\varepsilon$, and those
bounds give a primitive-recursive precision demand.  Thus this is
a positive-power monotonicity argument, not real trichotomy.

Here is one explicit demand that also explains the passage from
rational AM--GM to names.  Suppose all relevant values lie in
$[\ell,B]$ with positive rational $\ell$ and integer $B\ge1$.
Approximating each value to error $\eta\le\min(1,\ell/2)$ changes
a product of $s$ factors, and also an $s$th power of their mean,
by at most $s(B+1)^{s-1}\eta$, by telescoping the product.
For a target order tolerance $\varepsilon\le1$, it therefore
suffices to choose a dyadic $\eta$ with
\[
 \eta\le\min\left(1,\frac\ell2,
       \frac{\varepsilon^s}{8s(B+1)^{s-1}}\right).
\]
All quantities on the right are positive rationals.  A dyadic
satisfying the bound is found by a search bounded by their encoded
denominators.  This supplies a concrete primitive-recursive
precision demand for both finite products and the positive-power
comparison; strong positivity supplies $\ell$ for the finitely
many values used in this instance.
\end{proof}

\begin{theorem}[A local Lipschitz bound from two endpoint samples]
\label{thm:convex-local-modulus}
Assume the hypotheses of Lemma~\ref{lem:positive-convexity}.  Fix
positive rationals $\alpha\le\beta$.  Put
\[
 a=\frac\alpha2,\qquad b=\beta+1,\qquad
 \delta=\min\left(\frac\alpha2,1\right).
\]
Let $M\ge1$ be an integer upper bound for $F(a)$ and $F(b)$,
obtained from their precision-zero intervals.  Then, for rational
$x,y\in[\alpha,\beta]$,
\begin{equation}
 |F(x)-F(y)|\le\frac M\delta\,|x-y|.
 \label{eq:convex-lipschitz}
\end{equation}
In particular, with $H=\lceil M/\delta\rceil$,
\begin{equation}
 \mu_F(\alpha,\beta,k)=k+H
 \label{eq:explicit-bm-modulus}
\end{equation}
is a primitive-recursive local uniform-continuity modulus.
\end{theorem}

\begin{proof}
The case $x=y$ is immediate; rational order allows us otherwise to
assume $x<y$.  Express $y$ as the rational convex combination of
$x$ and $b$:
\[
 y=\frac{b-y}{b-x}x+\frac{y-x}{b-x}b.
\]
Lemma~\ref{lem:positive-convexity} gives
\begin{align*}
 F(y)-F(x)
 &\le\frac{y-x}{b-x}\bigl(F(b)-F(x)\bigr)\\
 &\le\frac{M}{\delta}(y-x).
\end{align*}
Here positivity gives $F(x)\ge0$, and $b-x\ge\delta$.
Likewise express $x$ as a rational convex combination of $a$ and
$y$:
\[
 x=\frac{y-x}{y-a}a+\frac{x-a}{y-a}y.
\]
The resulting inequality is
\[
 F(x)-F(y)
 \le\frac{y-x}{y-a}\bigl(F(a)-F(y)\bigr)
 \le\frac M\delta(y-x).
\]
Together these are \eqref{eq:convex-lipschitz}.  The denominators
are positive rationals with the displayed lower bound, so all
division is quantitatively justified.

For a Cauchy sampler $A_F$, rather than nested input intervals, one
may take
\[
 M=1+\left\lceil\max\bigl(|A_F(a,0)|,|A_F(b,0)|\bigr)\right\rceil.
\]
Coherence gives $F(a)\le A_F(a,0)+1\le M$ and the same for $b$.
Thus $M$ uses only two fixed-precision evaluations.
Since $H\ge1$ and $H\le2^H$, the hypothesis
$|x-y|\le2^{-(k+H)}$ gives
\[
 |F(x)-F(y)|\le H\,2^{-(k+H)}\le2^{-k}.
\]
This proves \eqref{eq:explicit-bm-modulus}.  All operations are
integer or rational arithmetic, a ceiling, and the two evaluations.
\end{proof}

This removes $M_F$ from the \emph{independent input data} for the
Bohr--Mollerup candidate on positive rationals.  Its strong positivity,
normalization, recurrence, and rational multiplicative convexity now
suffice to derive both its rational uniqueness theorem and the modulus
needed for its canonical extension to regular positive real names.
One can then use the repository's completion argument with the
derived modulus.

The scope of this repair is precise.  A rational presentation determines
its canonical continuous extension.  Assumptions imposed only at rational
arguments do not, by themselves, constrain arbitrary separately specified
values on irrational arguments.  If $F$ is presented on real names from
the outset, one must require that this presentation is the canonical
extension, or prove its compatibility with the rational presentation.
The theorem supplies the previously missing modulus in the stated
rational-presentation setting; it is not an extension theorem for an
unrestricted set-theoretic function.

\subsection{Why metastability does not supply a Cauchy rate}

There is a useful obstruction which can itself be stated with finite
arithmetic.  Let $h_e(n)$ be zero if machine $e$ has not halted in its
first $n$ steps, and one otherwise.  Bounded simulation is primitive
recursive, and $h_e$ is nondecreasing with at most one jump.

Fix a primitive-recursive control function $g$ with $g(n)\ge n$.
For any $e$, one of the two disjoint integer windows
\[
 [0,g(0)]\quad\hbox{and}\quad
 [g(0)+1,g(g(0)+1)]
\]
has no jump.  Indeed, each nonconstant window contains a jump strictly
between its first and last sampled index, and these windows cannot
both contain the unique jump.  Inspecting both finite windows therefore
finds an $N\in\{0,g(0)+1\}$ such that
\[
 |h_e(i)-h_e(j)|=0
 \quad\hbox{for all }i,j\in[N,g(N)].
\]
This is a uniform primitive-recursive metastability procedure for
every requested positive tolerance.

On the other hand, a uniform computable Cauchy rate $R(e)$ for tolerance
$1/2$ would decide halting: simulate $e$ for $R(e)$ steps.  If it has
not halted by then, it cannot halt later, because the values at
$R(e)$ and at the later halting time would differ by one.  Hence no
such computable rate, and in particular no primitive-recursive rate,
exists uniformly in $e$.  This is a metamathematical obstruction; it
does not add a machine-halting principle to \iPRA.

Consequently a finite fluctuation bound, or a theorem giving a good
finite window for every control function, cannot be silently substituted
for the modulus needed to construct a regular name.  In the positive
results above the modulus comes from an explicit continuity or
uniqueness inequality, and its dependence on the requested precision
is displayed.

\subsection{Arithmetic induction and the Lean implementation path}

For Theorem~\ref{thm:finite-net-infimum}, encode rational tuples and
the grid as finite natural-number codes.  The minimum loop maintains
the current smallest sampled rational, its first index, and the
Boolean conjunction of comparisons against all samples visited so
far.  The invariant is quantifier-free after replacing a bounded
conjunction by its primitive-recursive Boolean fold.  The running
maximum and minimum used for interval intersection have the same
kind of invariant.  Grid correctness is proved coordinate by coordinate
using the quotient/remainder identity for rational rounding.  The
continuity contract is instantiated at the one rounded point; one
never inducts on a formula asserting a property of all real points.

For Theorem~\ref{thm:coherent-minimizer}, the finite maximum defining
$v$ is another bounded fold.  Having fixed $k,l$, the proof calls the
uniqueness contract only once, at $\min(k,l)$.  The strong-convexity
corollary uses a midpoint, two value comparisons, and an inequality
of rational squares.  None of these steps uses induction on a
formula quantifying over all names or all precisions.

For Theorem~\ref{thm:convex-local-modulus}, rational convexity is a
theorem schema for each displayed finite pair of arguments and
rational weight.  The finite AM--GM arithmetic has loop bounds at
most $2s$ and finite product-error contracts.  Once ordinary
convexity is available, the modulus proof uses exactly two convexity
instances.  The choice of $M$ is a primitive-recursive expression
in two coarse samples; there is no search for an unknown neighborhood
or a positive real gap.

A useful Lean implementation order is the following.
\begin{enumerate}
\item Define rational codes, tuple lookup, $d_\infty$, dyadic grids,
      rounding, the first-index minimum fold, and running intersections.
      Prove their elementary finite specifications.  At the semantic
      level a structure may collect $A$, $\mu$, coherence and continuity,
      but the object-language theorem uses their fixed PR definitions.
\item Prove the \emph{finite} lemma
      \[
      m_k\le A(x,q)+\tfrac12\,2^{-k}+2^{-q}
      \]
      directly for arbitrary rational $x$ and precision $q$.  Prove
      $A(a_k,q)\le U_k+2^{-q}$ and the cross-interval inequality
      $L_i\le U_j$.  These are the mathematical core of the infimum
      constructor, and give its object-language targets without
      importing Lean's real compactness theorems.
\item Prove a generic rational lemma saying that two approximate
      minimizers under \eqref{eq:midpoint-drop} have squared distance
      at most $8\varepsilon/\lambda$.  Instantiate it with
      $\varepsilon=2^{-2k-c}$ to verify the point constructor.
\item Formalize weighted AM--GM, positive-power order reflection with
      an explicit precision demand, the two secant inequalities, and
      $H2^{-H}\le1$.  The candidate modulus then follows by composition.
\item Provide the arithmetic derivations of these fixed schemas and
      feed them to the common finite-contract compiler.  A Lean theorem
      about ordinary $\mathbb R$ is useful as a semantic check, but
      is not the object-language \iPRA proof.
\end{enumerate}

The contracts above identify the exact remaining formalization work.
They are definitions and proved mathematical specifications, rather
than unexplained assumptions named after the desired conclusion.

\section{Finite algebraic sections: a concrete repair of the Diophantine interface}
\label{sec:finite-sections}

The repository's general Siegel route already separates two assertions:
every supplied finite tuple has a cover by a bounded number of proper
linear subspaces, and each such subspace contains only a bounded number
of curve points. Its unresolved geometric input is a zero-count theorem
for functions with controlled poles. We now prove an elementary substitute
for a substantial explicit family. The resulting theorem supplies the
second assertion without a theory of divisors or Riemann--Roch spaces.
The quantitative approximation theorem supplying the first assertion
remains a separate, unproved input \cite{repo-siegel}.

All rational numbers in this section have canonical signed-numerator,
positive-denominator codes. Equality and order are decidable integer
tests. A polynomial is a finite coefficient array, with trailing zero
coefficients removed when its degree is computed. A point is a pair of
such rational codes. These choices make distinctness, substitution,
matrix rank, and list membership primitive recursive predicates.

\subsection{A finite polynomial zero-count lemma}

\begin{lemma}[Zeros in a supplied finite list]
\label{lem:section-root-count}
Let $h\in\Q[X]$ be nonzero of degree $D$. A finite list of distinct
rational roots of $h$ has length at most $D$. This statement has a proof
in $\iPRA$, uniformly in the coefficient and root lists.
\end{lemma}
\begin{proof}
We first give the division operation, including its bound. If
$h(X)=a_0+\cdots+a_DX^D$ with $D>0$, and $z\in\Q$, define
$q_{D-1}=a_D$ and, going downwards, $q_{j-1}=a_j+zq_j$ for
$1\le j<D$. Put $q(X)=\sum_{j<D}q_jX^j$. Comparing coefficients gives
\[
 h(X)=(X-z)q(X)+a_0+zq_0,
 \qquad a_0+zq_0=h(z).
\]
The computation uses exactly $D$ coefficient updates. If $h(z)=0$,
the remainder vanishes and $q$ has degree $D-1$ with the same nonzero
leading coefficient. If $w\ne z$ is another root, evaluation gives
$0=(w-z)q(w)$; cancellation in $\Q$ gives $q(w)=0$.

Suppose a supplied root list has at least $D+1$ entries, and retain its
first $D+1$ entries $z_0,\ldots,z_D$. Divide successively by
$X-z_0,\ldots,X-z_{D-1}$. After $j$ divisions the quotient has degree
$D-j$, has the original nonzero leading coefficient, and vanishes at
each of the remaining listed roots. At $j=D$ it is a nonzero constant
which vanishes at $z_D$, a contradiction. If $D=0$, the first root
already gives this contradiction.

This is finite induction, not induction on an assertion about all
rational roots. The remaining-root condition is a bounded conjunction
over the fixed input list. Its characteristic Boolean, the polynomial
coefficient identity, and the leading-coefficient equality form a
quantifier-free invariant. All divisions just described are by a monic
linear factor and require no search. The proof therefore uses only the
allowed induction instances.
\end{proof}

\subsection{An explicit annihilator of a proper span}

We describe the required linear algebra so that the word ``proper''
does not conceal an infinite-dimensional selection principle. For a
$t$ by $n$ rational matrix $W$, define its rank to be the largest size
of a nonzero minor, with rank zero when every entry is zero. This is
a bounded finite test. A proper span in $\Q^n$ can be represented by
such a matrix with rank less than $n$.

\begin{lemma}[A bounded, fraction-free annihilator]
\label{lem:section-annihilator}
For $n\ge1$, from a rational $t$ by $n$ matrix $W$ of rank less than
$n$, a primitive recursive operation constructs a nonzero integer
vector $\lambda\in\mathbb Z^n$ satisfying $W\lambda=0$.
Its correctness has a quantifier-free finite-arithmetic proof.
\end{lemma}
\begin{proof}
Clear all denominators using their positive product $D_0$ and put
$V=D_0W$. This does not change which minors vanish. By bounded
enumeration choose a nonzero minor of maximum size $k<n$, with row
set $I$ and column set $J$; let its matrix be $A$. For $k=0$ use the
empty minor, whose determinant is one. Choose the least column
$c\notin J$, and let $b$ be column $c$ restricted to the rows in $I$.
Define
\[
 \lambda_c=\det A,\qquad
 \lambda_J=-\operatorname{adj}(A)b,\qquad
 \lambda_j=0\quad(j\notin J\cup\{c\}).
\]
When $k=0$ this simply means $\lambda_c=1$ and all other entries zero.
In every case $\lambda$ is nonzero.

The selected rows annihilate $\lambda$ because cofactor expansion gives
$A\operatorname{adj}(A)=(\det A)I_k$. For any other row $v$ of $V$,
the block-determinant identity is
\[
 \det\begin{pmatrix} A&b\\v_J&v_c\end{pmatrix}
       =v_c\det A-v_J\operatorname{adj}(A)b=v\lambda.
\]
This identity follows by expansion along the last row; equivalently,
multiply the usual row-subtraction computation by $\det A$ and use
the adjugate identity. The displayed determinant is a minor of size
$k+1$, up to row and column ordering. It vanishes by maximality of
$k$. Thus every row annihilates $\lambda$. For $k=0$, maximality says
that all matrix entries are zero, which gives the same conclusion.
Since $D_0>0$, $V\lambda=0$ implies $W\lambda=0$.

There are at most $2^{t+n}$ row-mask/column-mask pairs to inspect.
Each determinant is the signed sum over permutations, enumerated by
a bounded loop (even a scan of all tuples below $n^n$, retaining the
permutations, suffices). The cofactor identities follow by finite
expansion: terms with repeated columns cancel in pairs, and the
remaining terms are exactly the determinant terms. The block identity
is the same finite expansion. All matrix entries and identities are
finite arrays of rational or integer codes.

For a numerical size bound, if all input numerators in absolute value
and denominators are at most $H\ge1$, then
$D_0\le H^{tn}$ and every entry of $V$ has absolute value at most
$C=H^{tn+1}$. Every coordinate constructed above has absolute value
at most $n!C^n$: a $k$-minor has bound $k!C^k$, and the adjugate
product has that same bound. This is a uniform primitive recursive
majorant. No arbitrary rational kernel vector is selected.
\end{proof}

\subsection{A zero-count theorem for superelliptic curves}

Fix integers $r\ge1$ and $s\ge0$, and a nonzero rational polynomial
$f$ of degree $m$ satisfying
\[
                         m>rs.
\]
Consider the affine point predicate and feature vector
\[
 \mathcal C_{f,r}(x,y)\;:\Longleftrightarrow\;y^r=f(x),
 \qquad
 \Phi_s(x,y)=(1,x,\ldots,x^s,y)\in\Q^{s+2}.
\]
No smoothness, irreducibility, or genus hypothesis is needed for the
following finite assertion.

\begin{theorem}[Uniform finite section bound]
\label{thm:superelliptic-section}
For every nonzero rational vector
$\lambda=(\lambda_0,\ldots,\lambda_s,a)$, a finite list of distinct
points satisfying
\[
             y^r=f(x),\qquad \lambda\cdot\Phi_s(x,y)=0
\]
has at most $m$ entries. Consequently the same bound holds for points
whose feature vectors lie in any supplied proper rational linear
subspace. Both statements have proofs in $\iPRA$.
\end{theorem}
\begin{proof}
Put $P(X)=\sum_{j=0}^s\lambda_jX^j$. The second equation is
$P(x)+ay=0$. The rational test $a=0$ gives two decidable cases.

If $a\ne0$, every such point has $y=-P(x)/a$. Substituting this
into the first equation gives
\[
             H(x)=0,\qquad H(X)=(-P(X))^r-a^rf(X).
\]
The first summand has degree at most $rs<m$, whereas the second has
degree exactly $m$ and nonzero leading coefficient. Hence $H$ is
nonzero of degree exactly $m$. Distinct points in this branch have
distinct $x$ coordinates, because $x$ determines $y$. Applying
Lemma~\ref{lem:section-root-count} to those coordinates gives the
bound $m$.

If $a=0$, the vector's nonzeroness says that $P$ is a nonzero
polynomial of degree at most $s$. Group the input points by equal
$x$ coordinates, assigning each point to the first earlier equal
coordinate, or starting a new group. There are at most $s$ groups,
by Lemma~\ref{lem:section-root-count}. In each group, the distinct
$y$ coordinates are roots of $Y^r-f(x)$, a polynomial of degree
exactly $r$. Thus each group has at most $r$ entries. The entire
list has at most $rs<m$ entries. This also handles $s=0$: a
nonzero constant $P$ has no roots, so there are no groups.

For a proper subspace supplied by a spanning matrix of rank less
than $s+2$, Lemma~\ref{lem:section-annihilator} supplies a nonzero
annihilator. Every vector in its span satisfies the resulting
linear equation. Thus the already proved hyperplane bound applies.

The polynomial powers use finite convolution of coefficient lists.
The degree statements are leading-coefficient statements; the
condition $m>rs$ makes the relevant coefficient survive without
cancellation. Grouping a list of length $t$ requires at most $t^2$
equality tests. The group-size assertion is a bounded conjunction,
and its sum is a primitive recursive fold. Together with the two
lemmas, this gives an explicit reduction of the whole proof to
quantifier-free induction on bounded arrays and loops.
\end{proof}

\begin{remark}[The constant is sharp]
For any $r,s,m$ with $m>rs$, choose a rational polynomial $Q$ of
degree at most $s$ and set
\[
          f(X)=Q(X)^r+\prod_{j=1}^{m}(X-j).
\]
It has degree $m$. The $m$ distinct points $(j,Q(j))$ lie on
$y^r=f(x)$ and on the feature hyperplane $y-Q(x)=0$. Thus a
smaller uniform bound cannot hold for this family as stated.
\end{remark}

\paragraph{The quintic instance.}
For any quintic $f$, take $r=2$, $s=2$. Then
\[
                  \Phi_2(x,y)=(1,x,x^2,y),
\]
and every proper subspace contains feature images of at most five
distinct affine rational points of $y^2=f(x)$. This includes the
monic squarefree quintic family used in the repository's Faltings
report \cite{repo-faltings}, but the proof does not need squarefreeness.
For a sharp numerical example take
\[
 f(X)=X^5-14X^4+85X^3-223X^2+274X-119
      =(X^2+1)^2+\prod_{j=1}^{5}(X-j).
\]
The points $(1,2),(2,5),(3,10),(4,17),(5,26)$ all lie in the
hyperplane $y-x^2-1=0$. The companion exact-arithmetic program
constructs an annihilator by minors and verifies the polynomial
identities for this example; the general proof is the argument above.

\subsection{Tuple-local covers and the remaining arithmetic theorem}

Let $\mathcal A$ be any additional decidable condition on the
points of $\mathcal C_{f,r}$. Examples might impose integrality,
congruence restrictions, or a particular approximation inequality.
We make no automatic assertion that these conditions yield a
small-subspace cover.

\begin{theorem}[The explicit finite cover interface]
\label{thm:superelliptic-cover}
Suppose that, for every finite list $L$ of distinct
$\mathcal C_{f,r}$-points satisfying $\mathcal A$, the feature
vectors $\Phi_s(L)$ admit a supplied cover by at most $M$ proper
rational linear subspaces. Then every such list has length at most
$mM$. It is enough to assume the cover assertion only for lists
of length $mM+1$. The covering subspaces may depend on the entire
input list; no global collection of exceptional subspaces is required.
\end{theorem}
\begin{proof}
For a supplied cover replace each subspace by the hyperplane from
Lemma~\ref{lem:section-annihilator}. Assign each point to the first
covering hyperplane by a bounded scan. The covering premise ensures
that an index is found. Each assigned sublist is distinct and
satisfies the corresponding hyperplane equation, so
Theorem~\ref{thm:superelliptic-section} bounds its length by $m$.
Summing over at most $M$ sublists gives $|L|\le mM$.
If the premise was stated only for length $mM+1$, any longer list
has a prefix of that length. The same argument would give
$mM+1\le mM$, a contradiction. The sum, assignment, and prefix
operations are bounded finite algorithms.
\end{proof}

The cover premise itself can be presented as a bounded arithmetic
predicate, without arbitrary subspace codes. For a list of $t$ feature
vectors, enumerate every subset of those vectors, compute its span,
and retain subsets of rank less than $s+2$. A span of vectors in a
proper supplied subspace remains proper: a nonzero annihilator of that
subspace also annihilates every linear combination of these vectors. A proper subspace
covering some of the tuple can be replaced by the span of the
points assigned to it. Thus one only needs $M$ subset masks, each
less than $2^t$, allowing the empty subset for padding. At most
$2^{tM}$ tuples of masks need be tested. All ranks and membership
conditions are decidable finite linear algebra. Equivalently,
replace each candidate span by its explicitly computed annihilator
hyperplane; this can only enlarge a covering set and remains proper.
This gives a fixed primitive recursive test
\[
       \operatorname{Cover}_{M,s}(L)\in\{0,1\}.
\]
The new theorem proves, in arithmetic,
\[
 \operatorname{ValidCurveList}_{f,r}(L)\ \wedge\
 \operatorname{Nodup}(L)\ \wedge\
 \operatorname{Cover}_{M,s}(L)=1
       \quad\Longrightarrow\quad |L|\le mM.
\]
It does not prove that the cover test succeeds for all lists with
$\mathcal A$. Proving that assertion is exactly where the requisite
quantitative approximation theorem must enter. This bounded
formulation identifies a precise interface for that future work.

If an explicitly listed exceptional set has at most $E$ points,
and the cover assertion applies outside it, the identical assignment
argument gives $E+mM$. For finitely many disjointly assigned
conditions with respective cover bounds $M_i$, the bound is
$E+m\sum_i M_i$. These are cardinality conclusions. They contain
no assertion locating all the exceptional points.

\subsection{Why a cardinality or gap bound still supplies no cutoff}

The repository correctly separates cardinality from a positive
height bound. Here is the exact general obstruction, including an
often tempting extra assumption. Let $T(e,n)$ mean that the fixed
universal machine with program/input code $e$ halts for the first
time at step $n$. Bounded simulation makes $T$ a fixed primitive
recursive predicate. For each $e$ it has at most one solution.
Consequently it satisfies every pairwise gap condition of the form
\[
 T(e,n)\wedge T(e,n')\wedge n<n'
             \quad\Longrightarrow\quad n'>G(e,n),
\]
for any proposed primitive recursive growth function $G$: there
are no two distinct solutions.

If a computable $B(e)$ bounded every solution of $T(e,n)$, simulate
the machine for $B(e)$ steps. This would decide whether it ever
halts. Such a decision procedure is impossible: a program which
halts exactly when that procedure predicts it will not halt,
applied to its own code, yields an immediate contradiction.
Therefore cardinality bounds, even together with arbitrarily strong
pairwise gap information, do not admit a general computable
conversion into height cutoffs. This argument uses one fixed bounded
simulator; it does not postulate a primitive recursive evaluator
for all primitive recursive functions.

Nothing here shows that a particular Diophantine cutoff is
unprovable. It shows why the finite-cover theorem must retain its
cardinality conclusion. Additional theorem-specific mathematics is
required to upgrade that conclusion
to a complete list or a uniform height bound.

\subsection{Lean definitions and proof obligations}

The following definitions suffice for a complete formalization path.
At the semantic level use \texttt{Rat}, finite coefficient arrays,
and finite lists of pairs. At the object level replace these by
the natural-number codes already described and use the same
bounded algorithms. The semantic and object-theoretic obligations
are distinct.

\begin{enumerate}
\item Define \texttt{QPoly} as a coefficient list, with
\texttt{coeff}, \texttt{eval}, \texttt{trim}, \texttt{degree},
\texttt{add}, \texttt{mul}, and \texttt{pow}. Zero is tested by a
Boolean fold over the coefficients. Implement
\texttt{syntheticDiv (p : QPoly) (z : Rat)} by the reverse coefficient
fold used in Lemma~\ref{lem:section-root-count}. Prove its coefficient
identity, the evaluation identity, and preservation of the leading
coefficient. Then prove \texttt{distinctRoots\_length\_le\_degree}
by iterating this operation over an input list, with the explicit
bounded invariant in that lemma.
\item Define rational matrices as dimensions plus flat arrays.
Implement minors, permutation determinants, bounded maximum-minor
rank, and \texttt{annihilator}. Prove determinant cofactor expansion
by finite sums, then the two identities used in
Lemma~\ref{lem:section-annihilator}. The target specification is
\texttt{rank W < n -> nonzero (annihilator W) = true}
together with \texttt{mulVec W (annihilator W) = zeroVec}.
The input $n>0$ is explicit.
\item Define
the curve, feature, and annihilator predicates by the finite operations
just described:
\begin{lstlisting}
onCurve f r (x,y) := (y^r == eval f x)
features s (x,y) := [1,x,...,x^s,y]
annihilates lam p := (dot lam (features s p) == 0)
\end{lstlisting}
The main specification quantifies over coefficient arrays and
finite lists:
if $r>0$, $\deg f=m>rs$, $\lambda$ has length $s+2$ and is
nonzero, and a duplicate-free list passes both Boolean tests, then
its length is at most $m$. Prove it by the decidable final-coefficient
split. In the nonzero case construct $(-P)^r-a^rf$; in the zero
case use a stable grouping operation with a proved partition-count
identity. No algebraic closure, real numbers, or abstract
finite-dimensional choice theorem is required.
\item Define \texttt{coverTest M s L} by the finite mask enumeration
above. Prove soundness by extracting the successful masks and their
annihilators. Prove completeness relative to supplied proper
subspaces by assigning each input point and replacing each group
by its span. Combine soundness with the section theorem to obtain
\texttt{coverTest\_implies\_length\_le}. A future quantitative
approximation development must separately prove that this Boolean
test is one on the required valid tuples.
\item Reify the rational arithmetic, coefficient folds, matrix
operations, and list grouping as particular primitive recursive
definitions. For every loop record its Boolean invariant and
compile its induction into a quantifier-free induction instance.
The human proofs above identify those instances. A compiled
semantic Lean theorem alone would not certify this final
object-theory translation.
\end{enumerate}

The root-count theorem, annihilator construction, superelliptic
section theorem, and finite-cover implication are proved
mathematically here. Their complete embedded Lean derivations are
formalization work. The uniform success of the quantitative
cover test is a further mathematical hypothesis, not part of that
formalization work and not claimed by this package.

\section{What the framework repairs, and what it leaves open}
\label{sec:audit}

\subsection{Repairs supported by an explicit new argument}

\paragraph{The FTA iteration restriction.}
The baseline gives a one-root procedure uniform in degree, but restricts
its deflation paragraph to an externally fixed degree
\cite{repo-fta}. Theorem~\ref{thm:fta-uniform-factorization} supplies the
missing uniform nested-evaluation argument. It has more content than a
renaming of the schema: the output root evaluator has the numerical
degree as an argument, its stage caps are computed in advance, and its
verification horizons are composed before the stage induction.

The zeta-zero report already constructs a different rational-polynomial
root backend after identifying that restriction \cite{repo-zeta-roots}.
Our result extends to the fixed Cauchy-coefficient schema. It does not
replace the rational backend's exact gcd, multiplicity, or separation
computations; those require different finite data.

\paragraph{The Bohr--Mollerup candidate modulus.}
The source asks whether its extra candidate modulus follows from the
other hypotheses \cite{repo-bm}. Theorem~\ref{thm:convex-local-modulus}
answers this for the source's rational presentation. It uses strong
positivity and multiplicative convexity, and in fact does not need the
normalization or recurrence. The modulus is a concrete expression in
rational endpoints and two coarse evaluations. The finite AM--GM padding
construction also answers the report's request for a simpler arithmetic
certificate architecture; we make no claim of a measured reduction in
checker bytes or native proof length.

\paragraph{The finite geometric section bound.}
The Siegel route separates an approximation/cover statement from a
zero-count or fibre bound \cite{repo-siegel}. Theorem~\ref{thm:superelliptic-section}
supplies that bound for $y^r=f(x)$, $\deg f>rs$, with explicit features.
For the quintic family in the Faltings report \cite{repo-faltings}, the
bound is five. A valid cover would yield $5M$ as a cardinality bound.
The existence of such covers for the relevant set of points is not
proved here. For some unrestricted point sets it is false; the geometric
bound must therefore not be presented as a theorem of general finiteness.

\paragraph{The common analytic interfaces.}
The calculus and Taylor--Cauchy reports already state fixed rational
samplers and modulus hypotheses \cite{repo-calculus,repo-taylor}.
The finite-net theorem makes their useful compact-domain extremum module
available once, with a precise output-value name. The coherent-minimizer
theorem separates the additional quantitative uniqueness information
needed to obtain a point. This is an extension of shared infrastructure,
not a claim that those source arguments were incorrect.

\subsection{Obstructions that cannot be repaired by changing the encoding alone}

\paragraph{Exact Puiseux data from arbitrary coefficient names.}
The Puiseux report gives a halting obstruction already in the family
$Y^2-X^2-a$, with $a$ supplied only by a Cauchy name
\cite{repo-puiseux}. Even the promise that the value happens to be a
dyadic rational does not provide its exact finite presentation. A zero
versus positive coefficient can alter the exact initial branch data.
The finite-observation framework retains exact algebraic presentations
or explicit discrete certificates when the requested output requires
them. Returning coherent approximate complex roots, as our factorization
theorem does, does not decide those branch data.

\paragraph{Apartness and inversion.}
A procedure dividing by a represented number needs quantitative
nonvanishing. The source FTA discussion uses $az-1$ to explain this
boundary \cite{repo-fta}. Proposition~\ref{prop:name-arithmetic} supplies
an explicit reciprocal construction from an apartness certificate,
while denominator-free identities and monic presentations can avoid some
uses of division. No universal nonzero-to-apartness principle is added.

\paragraph{Uniform evaluation of all PR program codes.}
The diagonal argument in Section~\ref{sec:compiler} rules out a primitive
recursive evaluator for every primitive recursive code. The
zeta-functional-equation plan already recognizes this issue
\cite{repo-zeta-plan}. Our theorem solves a variable-depth composition
problem for a fixed uniform table routine; it does not invalidate that
obstruction. The difference between a numerical degree and an arbitrary
program code is essential.

\paragraph{Metastability, cardinality, and cutoffs.}
The one-jump sequence in Section~\ref{sec:quantitative-analysis} has a
uniform finite-window procedure but no uniform computable Cauchy rate.
The singleton predicate in Section~\ref{sec:finite-sections} has a
uniform cardinality bound and vacuous pairwise gaps but no computable
uniform height cutoff. These prove failures of general conversion
principles. They do not establish that any named Diophantine theorem is
unprovable in \iPRA, nor that a specific sequence cannot have a modulus.

\subsection{Analytic number theory and probability: exposed dependencies}

The following are examples of mathematical input still needed after the
finite compiler is available. They are not an exhaustive classification
of the proof-theoretic strength of the corresponding classical theorems.

\begin{longtable}{@{}>{\raggedright\arraybackslash}p{0.25\linewidth}>{\raggedright\arraybackslash}p{0.43\linewidth}>{\raggedright\arraybackslash}p{0.26\linewidth}@{}}
\toprule
Reported difficulty & What a shared finite framework supplies & What it does not supply\\
\midrule
\endhead
Stronger prime number theorem error & Rational contour/error bookkeeping,
precision allocation, and common finite induction invariants. & A
stronger zero-free-region estimate. The inspected report establishes its
own weaker exponent, not the optimal requested rate \cite{repo-pnt}.\\[5pt]
Siegel--Walfisz and Bombieri--Vinogradov uniformity & Explicit contracts
for exceptional terms and their parameter dependence. & An effective
Siegel-type separation estimate, or the unproved exceptional-zero formula
itself \cite{repo-prime-sw}.\\[5pt]
Chebotarev estimates & Finite field-data and analytic-certificate
interfaces. & The required field-dependent analytic estimates and
exceptional-zero control \cite{repo-chebotarev}.\\[5pt]
General-law and advanced central limit statements & Finite-law
calculations can transfer through supplied PR tail, quantization, or
Lindeberg moduli. & High-dimensional convex smoothing, entropy-power
arguments, or a full stable-law classification \cite{repo-clt}.\\
\bottomrule
\end{longtable}

One particularly important uniformity issue is that excluding a single
exceptional modulus does not exclude all moduli divisible by its
conductor. The inspected prime-distribution report explicitly keeps this
problem visible \cite{repo-prime-sw}. Our framework preserves the
quantifiers and parameter dependence; it does not simplify the claim by
dropping the exceptional multiples.

\subsection{A concrete integration sequence}

A useful repository integration would proceed in three pieces. First
add the finite data, Boolean reflection, and table-compilation modules
under a shared arithmetic namespace; keep them independent of a single
analytic theorem. Second add the factorization, convex-modulus, and
superelliptic-section applications with their source-specific hypotheses.
Third refactor existing packets to call those modules, preserving any
remaining mathematical contracts.

The supplied archive includes an integration note, not a repository
mutation. Its priority is a reviewable mathematical and formalization
package. Nothing in this report asserts that the existing repository's
uncompiled outlines have thereby become completed object-theory proofs.

\section{Lean formalization: definitions, checked kernel, and deep targets}
\label{sec:formalization}

The formalization path has a useful small completed component and a much
larger explicitly specified endpoint. The completed component checks
finite-table locality and compilation. The endpoint is a family of
object-language \texttt{PRAAx} derivations. We do not identify the two.

\subsection{The checked finite-table kernel}

The file \texttt{lean/FiniteObservation.lean} imports \texttt{Std}, not
Mathlib, and was checked with Lean 4.32.0. Its central definitions are:
\begin{lstlisting}
abbrev Stream := Nat -> Nat
abbrev Table := List Nat

def PrefixEq (h : Nat) (a b : Stream) : Prop :=
  forall i, i <= h -> a i = b i

def tabulate (a : Stream) (h : Nat) : Table :=
  (List.range (h + 1)).map a

def demand (u : Nat -> Nat) : Nat -> Nat -> Nat
  | 0, k => k
  | n + 1, k => u (demand u n k)

def Local (u : Nat -> Nat)
    (t : Nat -> Stream -> Stream) : Prop :=
  forall s k a b,
    PrefixEq (u k) a b -> PrefixEq k (t s a) (t s b)

def ofBody (u : Nat -> Nat)
    (body : Nat -> Nat -> Table -> Nat)
    (s : Nat) (a : Stream) (p : Nat) : Nat :=
  body s p (tabulate a (u p))
\end{lstlisting}

The theorem \texttt{ofBody\_local} proves that a monotone numerical query
bound makes this table body local. The numerical budget at stage $s$ is
\[
 \mathtt{horizon}(u,N,k,s)=\mathtt{demand}(u,N-s,k).
\]
The executable \texttt{build} fixes $N,k$ and recurses on the stage,
storing a \texttt{List Nat}. The semantic \texttt{iter} recurses into
\texttt{Stream}; it is explicitly a reference specification, not a
claimed first-order implementation. The main checked theorem is
\begin{lstlisting}
theorem compiled_correct
    (ht : Local u t) (a : Stream) (N k : Nat) :
  lookup (build u t a N k N) k = iter t a N k
\end{lstlisting}
with the implicit parameters and types as in the file.

The inductive proof uses a stronger finite-prefix assertion at every
stage. At a successor, the budget identity
$h_s=u(h_{s+1})$ supplies exactly the input agreement needed by locality.
This is a real verification of the generic compilation mechanism, not a
theorem whose premise is the same as its conclusion.

The same file proves \texttt{allBelow\_iff}, the exact interpretation of
a finite Boolean conjunction; \texttt{finite\_safety}, a finite invariant
lemma; and \texttt{descent\_fuel}, a scalar strict-descent bound. The
latter is auxiliary infrastructure rather than a new Mordell--Weil
result. It states that a natural rank which decreases at every step for
$H$ steps satisfies $H\le\operatorname{rank}(0)$.

The axiom printouts for the four principal checked theorems contain
\texttt{propext} and \texttt{Quot.sound}. They contain neither
\texttt{Classical.choice} nor \texttt{sorryAx}. These are Lean's ambient
foundational dependencies; they are not axioms added to \iPRA. An axiom
audit alone would not establish the object-theory claim in either case.

\subsection{Type-checked mathematical specifications}

The separate file \texttt{lean/Targets.lean} defines the numerical and
rational interfaces for the main analytic and algebraic targets. Its
structures and predicates are specifications of proof obligations;
constructing a type or a structure with unfilled obligations is not a
proof that it is inhabited. No opaque theorem, new axiom, or
\texttt{sorry} is used to assert those conclusions. The file's comments
mark that status explicitly.

At this semantic specification layer it is convenient to use
\texttt{Rat}, lists, and Lean functions. Their object-language
counterparts must be natural codes and the fixed \texttt{PR} terms for
the algorithms in the manuscript. An arbitrary Lean function
\texttt{Nat -> Nat} carries no automatic evidence of primitive
recursiveness. That evidence is a separate deliverable at the next layer.

\subsection{The existing deep backend}

The inspected source files provide the following exact interfaces:
\begin{longtable}{@{}>{\raggedright\arraybackslash}p{0.35\linewidth}>{\raggedright\arraybackslash}p{0.59\linewidth}@{}}
\toprule
Object & Meaning\\
\midrule
\endhead
\texttt{PR k} & Syntax for $k$-ary primitive recursive functions,
with constructors \texttt{zero}, \texttt{succ}, \texttt{proj},
\texttt{comp}, \texttt{prec}.\\
\texttt{PTm n} & Terms in $n$ free numerical variables, with
\texttt{var} and \texttt{app}.\\
\texttt{Fm PTm n} & First-order formulas over these terms; formulas
are not Lean predicates in this layer.\\
\texttt{Prf PRAAx Gamma phi} & Intuitionistic derivability of the
encoded formula from the encoded context.\\
\texttt{PRA.qfInd} & Constructs an induction proof from an encoded
formula, a syntactic \texttt{QF} witness, and object-level base/step proofs.\\
\texttt{Prf.subst}, \texttt{Prf.weaken}, \texttt{Prf.trans} &
Capture-avoiding substitution, weakening, and simultaneous cut.\\
\texttt{pra\_iff\_native} & Transfers completed arithmetic proofs to
the native scalar calculus.\\
\bottomrule
\end{longtable}
These are read from the pinned sources \cite{repo-pr,repo-deduction,repo-native};
we have not rebuilt the full repository dependency closure in this run.

\subsection{The object-language plan, with no semantic shortcut}

\paragraph{Stage 1: numerical representations and PR definitions.}
Implement the pair, list, integer, and rational codes from
Section~\ref{sec:foundations}, including total default branches. For each
operation construct its actual syntax in \texttt{PR k}. A semantic
function interpreting these syntax trees can be defined in Lean to state
correctness, but it is a metatheoretic function. Do not postulate an
object-level PR evaluator for arbitrary \texttt{PR} syntax.

Define addition, truncated subtraction, products, powers, list lookup,
bounded folds, rational order, finite determinants, and polynomial
convolution through \texttt{comp} and \texttt{prec}. This is repetitive
syntax construction and can be generated, but the output terms and their
arithmetic defining-equation proofs must remain inspectable. The
universal object-level facts come from those proofs, not from running a
finite set of examples.

\paragraph{Stage 2: Boolean reflection as an object theorem.}
For a chosen finite predicate $B$, build a specific Boolean PR term
$\chi_B$. Encode $B\leftrightarrow\chi_B=1$ in
\texttt{Fm PTm n} and prove it using the fixed-formula construction of
Lemma~\ref{lem:reflection}. The careful two-direction proof of bounded
universal reflection avoids circularly using bounded induction before
establishing it. A syntactic QF checker certifies the formulas on which
\texttt{PRA.qfInd} is called.

\paragraph{Stage 3: finite safety and contracts.}
For each selected numerical invariant, form the exact formula
\[
 j\le H(\mathbf x)\longrightarrow
   I(\mathbf x,j,s_j(\mathbf x))=1.
\]
Its order test is a Boolean PR term, so the encoded formula is
quantifier-free. Supply the object derivations of the base and step,
then apply \texttt{PRA.qfInd}. Global name or continuity assumptions
may remain in the context. No rule infers the base or step merely from
an ambient Lean proof about standard natural numbers.

The generic compiler's local verification obligation is a single
arithmetic theorem with the stage, horizon, and metadata as free
numerical variables. This distinguishes a uniform theorem from a
metatheoretic family of proofs for separately fixed stages. Metadata
bounds must be monotone or replaced by bounded maxima, as in the
construction of $U_N$.

\paragraph{Stage 4: the table compiler.}
Construct PR syntax for the demand list, the initial observation table,
and the forward stage fold. Prove prefix locality for its fixed body,
then mirror the finite proof of \texttt{build\_correct}. At the deep
level its invariant is a Boolean conjunction over the finite rows,
not the semantic proposition about \texttt{Stream}s. Include the
initial-contract derivation. Include Cauchy regularity among the finite
checks whenever the output is supposed to be a regular name.

The resulting generic proof is a compiler \emph{for supplied arithmetic
proof packets}: numerical function definitions, formula definitions,
QF witnesses, and uniform local proof terms. It is not an unrestricted
reflection principle for arbitrary Lean mathematics.

\paragraph{Stage 5: factorization.}
The mathematical statements to encode are
Lemmas~\ref{lem:fta-descent-contract}, \ref{lem:fta-root-table}, and
\ref{lem:fta-deflation-table}. The imported descent algorithm requires
its own deep rational-arithmetic proof; a Lean section variable stating
its correctness is an explicit assumption until that work is done.
The additional tasks are dyadic rounding, the monic trajectory
invariant, normalized root regularity, deflation regularity, and the
four-term polynomial defect estimate. All function symbols and bounds
are displayed in Section~\ref{sec:fta-extension}.

After these local proofs, encode the cap and backward-horizon
recursions and apply the finite-table compiler. Prove the finite
coefficient telescoping identity and its $dWG^d2^{-t}$ bound. Fix $k$
before the stage induction and introduce its universal quantifier only
at the end. The final formula is exactly the rational coefficient
inequality \eqref{eq:fta-uniform-factorization}, together with the
regularity assertions, rather than an invocation of Lean's
\texttt{IsAlgClosed} for the ordinary complex field.

\paragraph{Stage 6: convexity and finite nets.}
Define dyadic grids as finite tuples, rational rounding, first-index
minimization, interval intersection, and the explicit schedules
$\mu(k+2)$, $v(k)$, and $2k+c$. Prove the finite sampled minimum
inequality first. Instantiate the continuity premise only at the rounded
point and the required finite precisions. The two endpoint secant
inequalities yield the Bohr--Mollerup modulus once finite AM--GM and
positive-power comparison have their arithmetic proofs.

The positivity witness is part of the original presentation. Rational
non-strict inequalities between names are verified by explicit finite
error bounds and Lemma~\ref{lem:vanishing-error}, not by assuming real
trichotomy. The canonical extension is constructed through the derived
modulus; separately assigned irrational values remain outside the claim.

\paragraph{Stage 7: finite algebraic sections.}
Define the rational polynomial and point-list predicates, synthetic
division, determinant enumeration, maximal-minor rank, and the explicit
annihilator. The finite polynomial root bound is an induction over a
supplied distinct root list. Then split on the decidable last coordinate
of the annihilator and prove the two branches of
Theorem~\ref{thm:superelliptic-section}. Implement subset-mask cover
verification and first-cover assignment. The cover theorem becomes an
arithmetic implication from a successful finite test. Its success on
all intended point tuples is a separate future mathematical theorem.

\paragraph{Stage 8: integration and proof audit.}
Only after the \texttt{Prf PRAAx} terms exist transfer them through
\texttt{pra\_iff\_native}. Record axiom printouts of the complete
endpoints and a build of their transitive dependencies. An integration
that leaves the desired conclusion in a structure field, an assumption,
an \texttt{opaque} declaration, or a \texttt{sorry} has not completed
this stage. The current archive explicitly stops before claiming it.

\subsection{Concrete proof targets and completion criteria}

\begin{longtable}{@{}>{\raggedright\arraybackslash}p{0.29\linewidth}>{\raggedright\arraybackslash}p{0.40\linewidth}>{\raggedright\arraybackslash}p{0.25\linewidth}@{}}
\toprule
Target & Required mathematical invariant & Current status\\
\midrule
\endhead
Prefix locality and table compilation & Equality on the scheduled
finite input prefixes implies equality of the output rows. & Ambient
Lean proof checked; deep coding pending.\\[4pt]
Finite Boolean reflection & Fold value one iff all bounded tests
succeed, with the noncircular induction proof. & Ambient finite-fold
lemma checked; general deep reflection pending.\\[4pt]
Uniform factorization & Good finite tables, degree/cap facts, local
coefficient defects, finite telescope. & Human proof from stated
Descent dependency; full Lean proof pending.\\[4pt]
Candidate continuity modulus & Finite AM--GM, witnessed positivity,
two rational secant inequalities. & Human proof; typed targets and
concrete arithmetic roadmap.\\[4pt]
Infimum and coherent optimizer & Finite minimum and cross-interval
checks; quantitative uniqueness at fixed precisions. & Human proof;
exact examples and typed targets.\\[4pt]
Superelliptic section bound & Quotient degree and remaining listed
roots; annihilator identity; finite grouping. & Human proof;
exact examples and typed targets.\\
\bottomrule
\end{longtable}

\section{Verification evidence and reproducibility}
\label{sec:verification}

\subsection{What was executed}

The checked Lean kernel was run with Lean 4.32.0. The package includes
its source and the actual compiler/axiom output. The target specification
file is separately type-checked; its acceptance verifies the definitions
and statement types, not the pending mathematical conclusions.
The build scripts use only Lean's standard library for these supplied
files. The deeper integration plan targets the repository's Mathlib-based
\texttt{ConciseIPRA} development.

All Python examples use exact integers and \texttt{fractions.Fraction}.
They do not use floating-point numerical agreement as evidence for a
universal claim. Three independent example programs were executed:
\begin{itemize}
\item \texttt{finite\_observations.py}: 588 comparisons with the nested
      reference evaluator, 147 prefix-compatibility comparisons, and
      eight exact defect-telescoping examples. The table runs visit
      11,424 observations in total. The same program checks the displayed
      FTA budget inequalities for degrees one through four.
\item \texttt{analysis\_exact\_checks.py}: 4,097 grid evaluations,
      81 cross-interval comparisons, 16 strong-convex point comparisons,
      2,794 AM--GM binary nodes, 3,240 rational Lipschitz comparisons,
      and 303 finite one-jump/metastability cases.
\item \texttt{finite\_curve\_sections.py}: two exact curve/section examples,
      including the sharp five-point quintic example and the
      zero-$y$-coefficient branch with two $x$-values and four points.
      It computes the annihilators by minors and checks the polynomial
      identities by rational arithmetic.
\end{itemize}

These are implementation and example checks. They are not proof checking
for the universal \iPRA{} assertions. The manuscript supplies the human
proofs, and the finite Lean kernel proves its own stated universal lemmas.
The full FTA descent at the enormous displayed schedules was not executed.

\subsection{The cost of a deliberately generous bound}

The factorization bounds are designed to establish primitive
recursiveness, not practical root-finding performance. The exact budget
program, with initial normalized coefficient cap two and target accuracy
$2^{-4}$, gives the following sizes. The final local precision is $4+c$;
``initial horizon bits'' is the binary length of the numerical index
through which the original coefficient table is requested.
\begin{center}
\begin{tabular}{@{}rrrr@{}}
\toprule
Degree & Precision shift $c$ & Initial horizon bits & Final cap bits\\
\midrule
1 & 39 & 15 & 14\\
2 & 307 & 40 & 117\\
3 & 2,917 & 86 & 966\\
4 & 29,109 & 163 & 8,310\\
\bottomrule
\end{tabular}
\end{center}
These values were computed exactly from the stated recurrences. They
are not measured runtimes and are not offered as competitive complexity
bounds. Improving them is compatible with, but not required by, the
logical objective of this report.

\subsection{Rebuilding the archive's outputs}

From the archive root:
\begin{lstlisting}
./build_pdf.sh
python3 code/finite_observations.py
python3 code/analysis_exact_checks.py
python3 code/finite_curve_sections.py
cd lean
./check.sh
\end{lstlisting}
The PDF build uses ordinary LaTeX packages and requires no network,
shell escape, or externally generated figures. The script has a local
format-generation fallback for a TeX installation whose format cache is
missing. The source is split by topic under \texttt{tex/}; rebuilding
requires those files as well as \texttt{report.tex}.

The provenance ledger pins every cited repository file to the inspected
commit and records its Git blob identifier. The theorem ledger records
which conclusions are human proofs, which reuse the finite descent
lemma, which remain conditional on a cover, and which have checked Lean
components. The checksum manifest covers the delivered files and allows
the archive to be compared after extraction.

\subsection{The remaining work is identified, not suppressed}

The large remaining formalization task is to construct the deep
arithmetic proof terms for the local mathematical contracts and the
numerical coding library, then instantiate the compiler. The remaining
mathematical task in the general Diophantine route is to prove the
quantitative cover assertion. The analytic-number-theory contracts listed
in Section~\ref{sec:audit} also remain substantive separate inputs.
None of these gaps is filled by the fact that an example program runs,
that a specification type checks, or that a generic conditional theorem
has been accepted by Lean.

\begin{thebibliography}{99}

\bibitem{repo-index} Vladimir Reshetnikov, \emph{ProveIt repository}: Delivered iPRA proof packets: inventory and validation status. Pinned revision \texttt{94000996e69a}. \href{https://github.com/VladimirReshetnikov/ProveIt/blob/94000996e69ab0c527393a1d02a4a037a4eb0493/Logic/PrimitiveRecursiveArithmetic/docs/proofs/README.md}{Inspected source}.

\bibitem{repo-pr} Vladimir Reshetnikov, \emph{ProveIt repository}: The deeply embedded arithmetic theory. Pinned revision \texttt{94000996e69a}. \href{https://github.com/VladimirReshetnikov/ProveIt/blob/94000996e69ab0c527393a1d02a4a037a4eb0493/Logic/PrimitiveRecursiveArithmetic/ConciseIPRA/Lean/ConciseIPRA/PR.lean}{Inspected source}.

\bibitem{repo-deduction} Vladimir Reshetnikov, \emph{ProveIt repository}: Intuitionistic natural deduction and structural rules. Pinned revision \texttt{94000996e69a}. \href{https://github.com/VladimirReshetnikov/ProveIt/blob/94000996e69ab0c527393a1d02a4a037a4eb0493/Logic/PrimitiveRecursiveArithmetic/ConciseIPRA/Lean/ConciseIPRA/Deduction.lean}{Inspected source}.

\bibitem{repo-native} Vladimir Reshetnikov, \emph{ProveIt repository}: Bi-interpretability with the native scalar core. Pinned revision \texttt{94000996e69a}. \href{https://github.com/VladimirReshetnikov/ProveIt/blob/94000996e69ab0c527393a1d02a4a037a4eb0493/Logic/PrimitiveRecursiveArithmetic/ConciseIPRA/Lean/ConciseIPRA/NativeBiInterpretation.lean}{Inspected source}.

\bibitem{repo-fta} Vladimir Reshetnikov, \emph{ProveIt repository}: Fundamental theorem of algebra: finite rational descent and fixed-degree factorization. Pinned revision \texttt{94000996e69a}. \href{https://github.com/VladimirReshetnikov/ProveIt/blob/94000996e69ab0c527393a1d02a4a037a4eb0493/Logic/PrimitiveRecursiveArithmetic/docs/proofs/algebra/FTA_iPRA/FTA_iPRA.tex}{Inspected source}.

\bibitem{repo-bm} Vladimir Reshetnikov, \emph{ProveIt repository}: Constructive Bohr--Mollerup theorem: candidate modulus and further questions. Pinned revision \texttt{94000996e69a}. \href{https://github.com/VladimirReshetnikov/ProveIt/blob/94000996e69ab0c527393a1d02a4a037a4eb0493/Logic/PrimitiveRecursiveArithmetic/docs/proofs/analysis/bohr_mollerup_ipra_2026-11-15/paper/bohr_mollerup_ipra.tex}{Inspected source}.

\bibitem{repo-calculus} Vladimir Reshetnikov, \emph{ProveIt repository}: Constructive calculus in iPRA. Pinned revision \texttt{94000996e69a}. \href{https://github.com/VladimirReshetnikov/ProveIt/blob/94000996e69ab0c527393a1d02a4a037a4eb0493/Logic/PrimitiveRecursiveArithmetic/docs/proofs/analysis/calculus_ipra/report.tex}{Inspected source}.

\bibitem{repo-taylor} Vladimir Reshetnikov, \emph{ProveIt repository}: Taylor and Cauchy estimates in iPRA. Pinned revision \texttt{94000996e69a}. \href{https://github.com/VladimirReshetnikov/ProveIt/blob/94000996e69ab0c527393a1d02a4a037a4eb0493/Logic/PrimitiveRecursiveArithmetic/docs/proofs/analysis/taylor_cauchy_ipra/report.tex}{Inspected source}.

\bibitem{repo-siegel} Vladimir Reshetnikov, \emph{ProveIt repository}: Siegel route: finite covers and the unresolved geometric input. Pinned revision \texttt{94000996e69a}. \href{https://github.com/VladimirReshetnikov/ProveIt/blob/94000996e69ab0c527393a1d02a4a037a4eb0493/Logic/PrimitiveRecursiveArithmetic/docs/proofs/diophantine-number-theory/siegel_ipra/paper/05_general_route.tex}{Inspected source}.

\bibitem{repo-faltings} Vladimir Reshetnikov, \emph{ProveIt repository}: Faltings report: explicit quintic arithmetic target. Pinned revision \texttt{94000996e69a}. \href{https://github.com/VladimirReshetnikov/ProveIt/blob/94000996e69ab0c527393a1d02a4a037a4eb0493/Logic/PrimitiveRecursiveArithmetic/docs/proofs/diophantine-number-theory/faltings_ipra/faltings_ipra.tex}{Inspected source}.

\bibitem{repo-puiseux} Vladimir Reshetnikov, \emph{ProveIt repository}: Puiseux report: exact presentations and the Cauchy-name obstruction. Pinned revision \texttt{94000996e69a}. \href{https://github.com/VladimirReshetnikov/ProveIt/blob/94000996e69ab0c527393a1d02a4a037a4eb0493/Logic/PrimitiveRecursiveArithmetic/docs/proofs/algebra/puiseux_ipra/puiseux_ipra.tex}{Inspected source}.

\bibitem{repo-zeta-plan} Vladimir Reshetnikov, \emph{ProveIt repository}: Zeta functional equation: formalization plan. Pinned revision \texttt{94000996e69a}. \href{https://github.com/VladimirReshetnikov/ProveIt/blob/94000996e69ab0c527393a1d02a4a037a4eb0493/Logic/PrimitiveRecursiveArithmetic/docs/proofs/analytic-number-theory/zeta_functional_equation_ipra/formalization_plan.md}{Inspected source}.

\bibitem{repo-zeta-roots} Vladimir Reshetnikov, \emph{ProveIt repository}: Zeta zeros: rational polynomial root backend. Pinned revision \texttt{94000996e69a}. \href{https://github.com/VladimirReshetnikov/ProveIt/blob/94000996e69ab0c527393a1d02a4a037a4eb0493/Logic/PrimitiveRecursiveArithmetic/docs/proofs/analytic-number-theory/zeta_zeros_ipra_2026-11-15/sections/polynomial_algebra.tex}{Inspected source}.

\bibitem{repo-pnt} Vladimir Reshetnikov, \emph{ProveIt repository}: Logarithmic-integral PNT estimates: analytic argument. Pinned revision \texttt{94000996e69a}. \href{https://github.com/VladimirReshetnikov/ProveIt/blob/94000996e69ab0c527393a1d02a4a037a4eb0493/Logic/PrimitiveRecursiveArithmetic/docs/proofs/analytic-number-theory/pnt_li_ipra_2026-11-15/tex/analytic.tex}{Inspected source}.

\bibitem{repo-prime-sw} Vladimir Reshetnikov, \emph{ProveIt repository}: Prime distribution: Siegel--Walfisz and Bombieri--Vinogradov obstacles. Pinned revision \texttt{94000996e69a}. \href{https://github.com/VladimirReshetnikov/ProveIt/blob/94000996e69ab0c527393a1d02a4a037a4eb0493/Logic/PrimitiveRecursiveArithmetic/docs/proofs/analytic-number-theory/prime_distribution_ipra/paper/sections/07_sw_bv.tex}{Inspected source}.

\bibitem{repo-chebotarev} Vladimir Reshetnikov, \emph{ProveIt repository}: Prime distribution: Chebotarev dependencies. Pinned revision \texttt{94000996e69a}. \href{https://github.com/VladimirReshetnikov/ProveIt/blob/94000996e69ab0c527393a1d02a4a037a4eb0493/Logic/PrimitiveRecursiveArithmetic/docs/proofs/analytic-number-theory/prime_distribution_ipra/paper/sections/06_chebotarev.tex}{Inspected source}.

\bibitem{repo-clt} Vladimir Reshetnikov, \emph{ProveIt repository}: Central-limit theorem packet: statement coverage ledger. Pinned revision \texttt{94000996e69a}. \href{https://github.com/VladimirReshetnikov/ProveIt/blob/94000996e69ab0c527393a1d02a4a037a4eb0493/Logic/PrimitiveRecursiveArithmetic/docs/proofs/probability/central_limit_ipra_2026-11-15/THEOREM_LEDGER.md}{Inspected source}.

\bibitem{kohlenbach-notes} Ulrich Kohlenbach. \emph{Proof Interpretations}. 1998. \href{https://www.brics.dk/LS/98/1/}{Primary source}.

\bibitem{kohlenbach-point} Ulrich Kohlenbach. \emph{Pointwise hereditary majorization and some applications}. Author-hosted manuscript. \href{https://www2.mathematik.tu-darmstadt.de/~kohlenbach/point.pdf}{Primary source}.

\bibitem{kohlenbach-uwkl} Ulrich Kohlenbach. \emph{On the Uniform Weak K\"onig's Lemma}. 1999. \href{https://www.brics.dk/RS/99/11/}{Primary source}.

\end{thebibliography}

\end{document}
